% !TeX program = xelatex
\documentclass[11pt,openany]{ctexbook}
\usepackage[a4paper,margin=2.45cm]{geometry}
\usepackage{amsmath,amssymb,amsfonts,bm,mathtools}
\usepackage{booktabs,array,longtable,tabularx}
\usepackage{amsthm}
\usepackage{tikz}
\usepackage{pgfplots}
\pgfplotsset{compat=1.18}
\usepackage{fancyhdr}
\usepackage{placeins}
\usepackage{makeidx}
\makeindex
\usepackage{enumitem}
\usepackage{hyperref}
\usepackage{xcolor}
\usepackage[most]{tcolorbox}
\usepackage{microtype}
\usepackage{needspace}

\hypersetup{
  colorlinks=true,
  linkcolor=blue,
  urlcolor=blue,
  citecolor=blue,
  pdftitle={计算物理学简明教程：从动力学约化到现代高阶离散},
  pdfauthor={}
}

\setcounter{tocdepth}{1} % 目录显示 part、chapter 和 section
\setlength{\parindent}{2em}
\setlength{\parskip}{0.35em}

\newcommand{\dd}{\,\mathrm d}
\newcommand{\ee}{\mathrm e}
\newcommand{\ii}{\mathrm i}
\newcommand{\R}{\mathbb R}
\newcommand{\Order}{\mathcal O}
\newcommand{\Lop}{\mathcal L}
\newcommand{\Nop}{\mathcal N}
\newcommand{\Ksp}{\mathcal K}
\newcommand{\avg}[1]{\left\langle #1\right\rangle}

\newtcolorbox{chaptersummary}{
  colback=blue!3,
  colframe=blue!45!black,
  title={本章总结},
  fonttitle=\bfseries,
  boxrule=0.6pt,
  arc=2mm,
  left=2mm,
  right=2mm,
  top=1mm,
  bottom=1mm
}

\newtcolorbox{authornote}{
  colback=yellow!8,
  colframe=orange!75!black,
  title={潜在认知跳跃},
  fonttitle=\bfseries,
  boxrule=0.8pt,
  arc=2mm,
  left=2mm,
  right=2mm,
  top=1mm,
  bottom=1mm
}

\newtcolorbox{askai}{
  colback=violet!3,
  colframe=violet!55!black,
  title={你应该问 AI},
  fonttitle=\bfseries,
  boxrule=0.7pt,
  arc=2mm,
  left=2mm,
  right=2mm,
  top=1mm,
  bottom=1mm
}

\newcommand{\studyroute}[3]{%
  \par\smallskip\noindent\textbf{展开与检索。}#1\par
  \noindent\textbf{关键词：}#2。\par
  \noindent\textbf{AI 展开：}#3\par
}


\raggedbottom
\setlist{nosep,topsep=3pt}
\pagestyle{fancy}
\fancyhf{}
\fancyhead[LE,RO]{\small\nouppercase{\leftmark}}
\fancyfoot[C]{\thepage}
\setlength{\headheight}{15pt}
\newcommand{\learninggoals}[1]{\paragraph{本章学习目标}#1}
\newcommand{\conceptbridge}[3]{%
\begin{authornote}
\textbf{潜在疑问：}#1\par
\textbf{作者补充：}#2\par
\textbf{核心概念联系：}#3
\end{authornote}}
\newcommand{\aiexercise}[2]{%
\begin{askai}
\textbf{AI 展开问题：}#1\par\smallskip
\textbf{进一步展开：}#2
\end{askai}}
\newcommand{\norm}[1]{\left\lVert#1\right\rVert}
\newcommand{\jump}[1]{\left[#1\right]}
\newcommand{\avgface}[1]{\left\{#1\right\}}
\newcommand{\Poly}{\mathbb P}
\newcommand{\tr}{\operatorname{tr}}
\newcommand{\diag}{\operatorname{diag}}
\numberwithin{equation}{chapter}
\newtheorem{theorem}{定理}[chapter]
\newtheorem{proposition}[theorem]{命题}
\theoremstyle{definition}
\newtheorem{definition}[theorem]{定义}

\title{从动力学约化到现代高阶离散\\[0.6em]\large From Dynamical Reduction to Modern High-Order Discretization}
\author{Harvey Tan\\[0.6em]\large 计算物理学简明教程\quad 第 1 卷}
\date{2026 年 11 月}
\hypersetup{pdfauthor={Harvey Tan}}
\begin{document}
\frontmatter
\maketitle
\chapter*{前言}
\addcontentsline{toc}{chapter}{前言}
计算物理在有限资源下表示物理状态、约化自由度、补充闭合关系，并构造有限维演化。每一步都需要说明保留的信息、引入的假设以及可验证的结构。本卷以守恒、稳定性、对称性和波动误差评价这些选择。

全书围绕四个问题组织：如何表示状态（representation），如何利用尺度和统计信息约化动力学（reduction），如何用物理假设补足未解变量的影响（closure），以及如何约束有限维残差（discretization）。模型选择决定计算对象，离散方法决定其近似表示与演化；结构保持贯穿这一过程。

第一部分从 Hamilton 动力学进入 Liouville、动理学与连续介质。第二部分从分部积分建立函数空间与弱形式，比较试探空间、测试空间及不同残差约束。第三部分推导 DG、FR、VCJH、Mortar 和时间推进。第四、五部分分别发展离散速度与物理约化两条路线。第六部分联系网格、方法选择和验证。

本卷面向具有本科高年级数学基础的读者，先备知识包括多元微积分、线性代数、ODE、PDE 与基础数值分析。关键定义、边界项和算法构造在正文中给出；一般迹定理、误差理论和研究专题由“展开与检索”与具体 AI 问题提供入口。每章的学习目标与概念联系用于定位当前步骤在知识链中的作用。

本卷在代数系统的形成处引向第 0 卷《数值线性代数与现代高性能计算》，由矩阵结构继续讨论分解、Krylov、预条件、多重网格和并行执行。第二卷则以 Maxwell 方程继续研究几何、约束传播与电磁离散。

\paragraph{符号约定}
$\rho_N$ 为归一化的多粒子概率密度，$f_s$ 为数密度型分布；$\rho$ 为质量密度，$\rho_c$ 为电荷密度。$u$ 为标量场，$U$ 为守恒变量或其离散系数，具体维数随模型给出。网格尺度为 $h$、多项式次数为 $p$、时间步长为 $\Delta t$。第 10--14 章一维界面统一取左迹 $u^-$、右迹 $u^+$，跳跃为 $[u]=u^+-u^-$，平均为 $\{u\}=(u^-+u^+)/2$。数值法向通量另含法向符号。碰撞时间、黏性应力与离散松弛时间按章节明确区分。

\tableofcontents
\mainmatter
\part{从微观动力学到连续物理模型}
\chapter{从多粒子动力学到动理学方程}\label{ch:1}

\begin{chaptersummary}
从粒子轨道走向可计算模型，需要逐步选择保留的信息。本章依次引入概率密度、边缘化与闭合：概率描述表征微观状态的不确定性，边缘积分减少显式坐标，闭合近似处理低阶分布演化中的高阶关联。这条约化路线构成后续连续介质与等离子体模型的共同起点。
\end{chaptersummary}
\learninggoals{说明轨道与概率描述的关系；区分边缘积分、关联层级与碰撞闭合；核对分布函数的归一化和守恒矩。}
\studyroute{这里略去 Liouville 方程到 BBGKY 层级的逐项积分，以及由双粒子关联构造碰撞算子的极限过程。}{Liouville equation; marginal distribution; BBGKY hierarchy; molecular chaos}{设两粒子受成对力作用，写出单粒子边缘分布的积分定义，并推导其方程中依赖双粒子分布的项。}

\section{Hamilton 动力学与相空间输运}\index{Liouville 方程}
设系统中有 \(N\) 个粒子，第 \(i\) 个粒子的位置和动量为
\[
\bm x_i,\qquad \bm p_i.
\]
整个系统的状态是 \(6N\) 维相空间中的一点：
\[
\Gamma=
(\bm x_1,\dots,\bm x_N,\bm p_1,\dots,\bm p_N).
\]

若 Hamiltonian 为 \(H\)，粒子运动满足
\[
\dot{\bm x}_i=\frac{\partial H}{\partial \bm p_i},
\qquad
\dot{\bm p}_i=-\frac{\partial H}{\partial \bm x_i}.
\]

当研究目标是宏观统计量且逐粒子计算成本过高时，可引入相空间概率密度
\[
\rho_N(\Gamma,t).
\]

\begin{authornote}
\textbf{潜在疑问：}
“为什么从确定性的 \(N\) 粒子轨道转向概率密度？”

\textbf{作者补充：}
许多工程问题关注密度、温度或平均流速，而非每个粒子的精确轨道。对于粒子数巨大的系统，逐一追踪所有自由度通常代价过高；若直接计算所有粒子对之间的相互作用，仅力的计算就有 $O(N^2)$ 的成本。概率描述将关注点转向可能状态及其统计规律，也允许描述对初始微观状态的不确定性。需要追踪特定粒子时，仍可使用粒子方法或欧拉--拉格朗日混合方法。

\textbf{核心概念联系：}
\[
\text{deterministic particle trajectory}
\longrightarrow
\text{phase-space probability description}
\]
\end{authornote}

概率守恒给出高维连续性方程
\[
\frac{\partial \rho_N}{\partial t}
+
\nabla_\Gamma\cdot(\rho_N\dot\Gamma)=0.
\]

\begin{askai}
AI 展开问题：
\begin{quote}
从普通三维连续性方程开始，逐步类比到 \(6N\) 维相空间，解释为什么概率密度 \(\rho_N(\Gamma,t)\) 满足
\[
\partial_t\rho_N+\nabla_\Gamma\cdot(\rho_N\dot\Gamma)=0.
\]
请明确相空间中“通量”的含义。
\end{quote}

进一步展开：
\begin{quote}
用一个一维位置加一维动量的 Hamilton 系统作为最小示例，说明概率云团如何被相空间速度场搬运。
\end{quote}
\end{askai}

Hamilton 系统满足
\[
\nabla_\Gamma\cdot\dot\Gamma=0,
\]
故有
\[
\frac{\partial \rho_N}{\partial t}
+
\dot\Gamma\cdot\nabla_\Gamma\rho_N=0,
\]
或
\[
\frac{\partial \rho_N}{\partial t}+\{\rho_N,H\}=0.
\]

刘维尔方程描述概率密度随微观动力学输运；Hamilton 流保持相空间体积。它仍然是 \(6N\) 维方程，后续边缘化才减少显式坐标的数量。

\begin{authornote}
\textbf{潜在疑问：}
“概率描述为何仍然具有高维性？边缘化保留和消去了哪些信息？”

\textbf{作者补充：}
概率密度 \(\rho_N\) 依赖整个系统的 \(6N\) 个相空间坐标，因此改用概率描述本身并未降低维数。即使降到单粒子分布 \(f(\bm x,\bm v,t)\)，直接张量网格离散仍需约 \(N_x^3N_v^3\) 个自由度：每个空间网格点都要存储一份三维速度分布。例如每个坐标方向取 100 点便有 \(10^{12}\) 个数，仅以双精度存储一个标量数组就约需 8 TB。实际代价还取决于维度约化、网格、碰撞模型及并行实现。

边缘化通过积分保留所选粒子的统计状态，将其余粒子的具体位置、速度及单粒子分布所不能表达的联合关联消去。相互作用仍通过这些关联影响保留变量的演化，因此需要闭合近似。三类操作各有作用：边缘化定义保留变量，闭合规定未解关联的影响，初值和边界条件规定模型的时空数据。

\textbf{核心概念联系：}
\[
\rho_N(\Gamma,t)
\longrightarrow
f(\bm x,\bm v,t)
\]
\end{authornote}

\section{边缘分布、归一化与关联层级}\index{边缘化}
对同种质量 $m$ 的粒子，将动量坐标改为速度坐标，令
\[
\varrho_N(\bm x_1,\bm v_1,\ldots,\bm x_N,\bm v_N,t)
=m^{3N}\rho_N(\bm x_1,\ldots,m\bm v_N,t),\qquad
\int\varrho_N\dd z_1\cdots\dd z_N=1,
\]
其中 $z_i=(\bm x_i,\bm v_i)$，右端 $\rho_N$ 的全部动量均取 $\bm p_i=m\bm v_i$。
概率型边缘和数密度型边缘分别为
\[
g_s(z_1,\ldots,z_s)=\int\varrho_N\dd z_{s+1}\cdots\dd z_N,\qquad
f_s=\frac{N!}{(N-s)!}g_s.
\]
于是 $\int g_s=1$、$\int f_1=N$、$\int f_2=N(N-1)$。以下 $f=f_1$。
边缘化的积分在每一固定时刻进行，保留变量的演化通过对 Liouville 方程积分获得。

定义单粒子分布函数
\[
f(\bm x,\bm v,t),
\]
使
\[
f(\bm x,\bm v,t)\dd\bm x\dd\bm v
\]
表示位置和速度落在相应微元中的粒子期望数。

从 \(N\) 粒子方程对其余变量积分，会出现
\[
f_1\longrightarrow f_2\longrightarrow f_3\longrightarrow\cdots,
\]
即 BBGKY 层级。单粒子方程依赖双粒子关联，双粒子方程又依赖三粒子关联。要得到封闭模型，必须对关联作假设。

\begin{authornote}
\textbf{潜在疑问：}
“为什么从边缘化会自然出现 \(f_1,f_2,f_3,\dots\) 这个层级，而不是直接得到一个单粒子方程？”

\textbf{作者补充：}
单粒子受到其他粒子的作用，其平均效果依赖于这些粒子与所选粒子同时处于不同状态的概率，也就是双粒子分布。对其余变量积分虽消去了它们的显式坐标，却没有消除相互作用；同理，双粒子分布的演化又依赖三粒子分布。要只求解低阶方程，就需要在适用条件下近似高阶关联，这一步称为闭合。

\textbf{核心概念联系：}
\[
\text{marginalization}
\longrightarrow
\text{BBGKY hierarchy}
\longrightarrow
\text{closure problem}
\]
\end{authornote}

\section{动理学方程与碰撞模型}\index{Boltzmann 方程}
一般动理学方程写成
\[
\frac{\partial f}{\partial t}
+\bm v\cdot\nabla_{\bm x}f
+\frac{\bm F}{m}\cdot\nabla_{\bm v}f
=C[f].
\]
左侧是相空间输运，右侧是碰撞或松弛效应。此处使用 $\nabla_{\bm v}\cdot\bm F=0$ 的力场，包括速度无关力和 Lorentz 力；一般速度依赖力应写成速度空间的散度形式。

若忽略离散碰撞，带电粒子满足 Vlasov 方程：
\[
\frac{\partial f_s}{\partial t}
+\bm v\cdot\nabla_{\bm x}f_s
+\frac{q_s}{m_s}
(\bm E+\bm v\times\bm B)\cdot\nabla_{\bm v}f_s=0.
\]
电荷密度和电流密度由
\[
\rho_c=\sum_s q_s\int f_s\dd\bm v,
\qquad
\bm J=\sum_s q_s\int \bm v f_s\dd\bm v
\]
给出，再反馈到 Maxwell 方程中。由此形成 Vlasov--Maxwell 自洽系统。

对于稀薄中性气体，把碰撞近似为瞬时二体碰撞，可得玻尔兹曼方程
\[
\frac{\partial f}{\partial t}
+\bm v\cdot\nabla_{\bm x}f
+\frac{\bm F}{m}\cdot\nabla_{\bm v}f
=Q(f,f),
\]
其中
\[
Q(f,f)
=
\int B\left(f'f_*'-ff_*\right)\dd\Omega\dd\bm v_*.
\]
对碰撞前粒子作分子混沌近似，并在局部碰撞对的归一化下写成：
\[
f_2(\bm v,\bm v_*)\approx f(\bm v)f(\bm v_*).
\]

对于前述有限 $N$ 数密度归一化，独立粒子的精确乘积还含 $(N-1)/N$ 因子；常用碰撞闭合在大粒子数极限吸收该因子。碰撞前无关联的假设只在相应稀薄极限和入射粒子统计中使用。

\begin{askai}
AI 展开问题：
\begin{quote}
用两个粒子的速度分布做一个最小例子，解释 \(f_2(\bm v,\bm v_*)\approx f(\bm v)f(\bm v_*)\) 意味着什么。请分别说明“无关联”在概率论上、物理碰撞上、以及 Boltzmann 方程封闭上各是什么意思。
\end{quote}

进一步展开：
\begin{quote}
给出一个 \(f_2\neq f f_*\) 的反例，并说明这种相关性会让单粒子 Boltzmann 方程少掉哪些信息。
\end{quote}
\end{askai}

完整碰撞积分计算昂贵，BGK 模型用单一松弛过程替代：
\[
C[f]=\frac{f^{\mathrm{eq}}-f}{\tau},
\]
其中
\[
f^{\mathrm{eq}}
=
n\left(\frac{m}{2\pi k_BT}\right)^{3/2}
\exp\left[
-\frac{m|\bm v-\bm u|^2}{2k_BT}
\right].
\]

对于等离子体中的长程 Coulomb 相互作用，粒子常经历大量小角度偏转。此时可得到 Landau 或 Fokker--Planck 型算子，把碰撞理解为速度空间中的漂移和扩散。

因此，碰撞项的不同假设对应一组模型层级：
\[
\text{无碰撞 Vlasov}
\longrightarrow
\text{弱碰撞 Vlasov--Fokker--Planck}
\longrightarrow
\text{强碰撞流体模型}.
\]

\chapter{从速度矩到连续介质方程}\label{ch:2}

\begin{chaptersummary}
上一章的分布函数保留完整速度信息。本章通过速度矩定义密度、平均速度、压力和能量，推导宏观守恒方程，再以本构关系补足未解应力与热流。这一过程为 Euler、Navier--Stokes 和有限体积守恒形式提供基础。宏观模型的适用范围由保留的矩和闭合假设共同决定。
\end{chaptersummary}
\learninggoals{推导守恒矩并辨认未闭合项；区分压力张量、黏性应力和热流；给出 Euler 与 Navier--Stokes 的适用假设。}
\studyroute{这里略去能量矩的完整代数推导与输运系数的动理学计算。}{velocity moments; Chapman--Enskog expansion; constitutive closure}{从 $f$ 的前三类守恒矩出发，区分精确矩恒等式与 Newton 黏性、Fourier 热流所增加的假设。}

\section{守恒矩与宏观观测量}\index{矩闭合}
本章先取三维、单组分、单原子稀薄气体，弹性碰撞守恒质量、动量与能量，外力为 $m\bm b(\bm x,t)$。速度空间积分假设 $f$ 足够衰减。
分布函数包含完整速度分布，宏观模型选择其中有限个加权积分。

数密度、质量密度和宏观速度分别为
\[
n=\int f\dd\bm v,
\qquad
\rho=mn,
\qquad
\rho\bm u=m\int \bm v f\dd\bm v.
\]

定义相对速度
\[
\bm c=\bm v-\bm u.
\]
压力张量与热流为
\[
\bm P=m\int \bm c\bm c f\dd\bm v,
\qquad
\bm q=\frac{m}{2}\int |\bm c|^2\bm c f\dd\bm v.
\]

对动理学方程取零阶矩，利用碰撞守恒粒子数：
\[
\int Q(f,f)\dd\bm v=0,
\]

\begin{authornote}
\textbf{潜在疑问：}
“为什么宏观方程可以通过对速度变量取矩得到？取零阶矩、一阶矩的物理含义是什么？”

\textbf{作者补充：}
密度、平均速度和能量是完整速度分布的加权积分。这些矩直接连接守恒律、实验测量与宏观状态，因此成为连续介质模型的基本变量。有限矩也规定了描述的分辨能力：具有相同低阶矩的不同速度分布，在这一层次上对应相同的宏观状态。

\textbf{核心概念联系：}
\[
f(\bm x,\bm v,t)
\longrightarrow
\text{moments}
\longrightarrow
\text{macroscopic observables}
\]
\end{authornote}

得到连续性方程
\[
\boxed{
\frac{\partial\rho}{\partial t}
+\nabla\cdot(\rho\bm u)=0
}.
\]

其控制体形式为
\[
\frac{\dd}{\dd t}\int_\Omega \rho\dd V
=
-\int_{\partial\Omega}\rho\bm u\cdot\bm n\dd S.
\]
有限体积法正是从这一形式出发。

取一阶矩得到动量方程
\[
\frac{\partial(\rho\bm u)}{\partial t}
+
\nabla\cdot(\rho\bm u\bm u+\bm P)
=
\rho\bm b.
\]

\begin{askai}
AI 展开问题：
\begin{quote}
从一般动理学方程出发，逐步对速度变量取一阶矩，推导
\[
\partial_t(\rho\bm u)+\nabla\cdot(\rho\bm u\bm u+\bm P)=\rho\bm b.
\]
请明确说明每一项的积分分部、边界衰减假设和物理含义。
\end{quote}

进一步展开：
\begin{quote}
把 \(\bm P=m\int \bm c\bm c f\dd\bm v\) 拆成各向同性压力和偏应力，说明它为什么不是简单标量。
\end{quote}
\end{askai}

写成
\[
\bm P=p\bm I-\bm\tau,
\]
则
\[
\frac{\partial(\rho\bm u)}{\partial t}
+
\nabla\cdot(\rho\bm u\bm u)
=
-\nabla p+\nabla\cdot\bm\tau+\rho\bm b.
\]

\section{本构闭合与连续介质模型}
要封闭方程，必须给出应力与宏观变量之间的本构关系。Newton 流体满足
\[
\bm\tau
=
\mu\left[\nabla\bm u+(\nabla\bm u)^T\right]
+\lambda(\nabla\cdot\bm u)\bm I.
\]

\begin{authornote}
\textbf{潜在疑问：}
“为什么推到动量方程后还没有结束？本构关系为什么不是从前面的矩操作自动得到的？”

\textbf{作者补充：}
低阶矩的演化方程通常含有更高阶矩，例如动量方程中的应力和能量方程中的热流。若只求解有限个宏观变量，就必须用这些变量及其梯度近似尚未求解的量。本构关系正是这种额外的物理假设，其形式和适用范围取决于介质及所研究的尺度。

\textbf{核心概念联系：}
\[
\text{moment equations}
\longrightarrow
\text{unclosed stresses and heat fluxes}
\longrightarrow
\text{constitutive closure}
\]
\end{authornote}


不可压缩、常密度、常黏度情况下，
\[
\nabla\cdot\bm u=0,
\]
\[
\rho\left(
\frac{\partial\bm u}{\partial t}
+\bm u\cdot\nabla\bm u
\right)
=
-\nabla p+\mu\nabla^2\bm u+\rho\bm b.
\]

其中惯性项负责搬运动量，压力项负责重新分配动量，黏性项负责平滑速度梯度。若忽略黏性和热传导，则得到 Euler 方程。

取能量矩可得总能量方程
\[
\frac{\partial(\rho E)}{\partial t}
+
\nabla\cdot[(\rho E+p)\bm u]
=
\nabla\cdot(\bm\tau\cdot\bm u)
-\nabla\cdot\bm q
+\rho\bm b\cdot\bm u,
\]
其中 $e$ 为单位质量内能，单原子理想气体满足 $p=\rho k_BT/m$、$e=3k_BT/(2m)$。应力和热流的 Newton--Fourier 形式对应近局部平衡、短平均自由程尺度。
\[
E=e+\frac12|\bm u|^2,
\qquad
\bm q=-k\nabla T.
\]

从动理学到连续介质的本质，是只保留有限个矩，并用本构关系封闭更高阶矩。

\chapter{电磁场、Maxwell 方程与自洽耦合}\label{ch:3}

\begin{chaptersummary}
前两章从粒子与分布函数构造物质模型，本章进一步引入与物质自洽耦合的电磁场。Maxwell 约束联系电荷守恒，场粒子反馈形成非线性系统，线性化则提取基态附近的小扰动算子。由此得到的局部响应为色散关系、稳定性分析及第二卷的电磁离散奠定基础。
\end{chaptersummary}
\learninggoals{连接电荷连续性与 Maxwell 约束；列出 Vlasov--Maxwell 的自洽数据；在给定基态处构造 Fr\'echet 导数并解释平面波分析的条件。}
\studyroute{这里略去材料模型与动理学响应的具体选取，集中于线性化算子的共同结构。}{Fr\'echet derivative; linear response; Vlasov linearization}{对 $\bm D=\varepsilon\bm E+\alpha|\bm E|^2\bm E$ 求切线张量，说明基态方向如何进入线性化方程。}

\section{电磁守恒与场粒子耦合}\index{Vlasov--Maxwell 系统}
Maxwell 方程组为
\[
\nabla\cdot\bm E=\frac{\rho_c}{\varepsilon_0},
\qquad
\nabla\cdot\bm B=0,
\]
\[
\frac{\partial\bm B}{\partial t}+\nabla\times\bm E=0,
\]
\[
\frac{\partial\bm E}{\partial t}
-c^2\nabla\times\bm B
=
-\frac{\bm J}{\varepsilon_0}.
\]

对 Amp\`ere--Maxwell 方程取散度，并使用 Gauss 定律，得到
\[
\frac{\partial\rho_c}{\partial t}
+\nabla\cdot\bm J=0.
\]

\begin{askai}
AI 展开问题：
\begin{quote}
从 Amp\`ere--Maxwell 方程
\[
\partial_t\bm E-c^2\nabla\times\bm B=-\bm J/\varepsilon_0
\]
出发，逐步取散度并使用 Gauss 定律，推导电荷守恒方程。请特别解释为什么 \(\nabla\cdot(\nabla\times\bm B)=0\) 是关键结构。
\end{quote}

进一步展开：
\begin{quote}
给出一个离散算子例子，说明如果离散的 div-curl 结构不满足 \(D_hC_h=0\)，可能怎样制造非物理电荷。
\end{quote}
\end{askai}

因此 Maxwell 方程的连续结构自动包含电荷守恒。数值离散若破坏这一结构，可能产生非物理电荷。

带电连续介质中的 Lorentz 力为
\[
\bm f_L=\rho_c\bm E+\bm J\times\bm B.
\]
结合流体方程和 Ohm 定律，可得到双流体模型或单流体 MHD。
对动理学物种，$f_s$ 决定 $\rho_c$、$\bm J$，场又通过 $q_s(\bm E+\bm v\times\bm B)$ 改变 $f_s$。这一双向依赖构成场粒子自洽耦合。单流体 MHD 还需低频、准中性及尺度排序，将粒子模型与物种流体方程约化为共同速度、压力和磁场；具体封闭见第~\ref{ch:18}~章。

\section{非线性算子与局部线性化}
真空 Maxwell 方程是线性的，但 Maxwell 方程与介质、载流子、流体和非线性本构关系耦合后，整体系统通常是非线性的。设耦合模型抽象写成
\[
\Nop(\bm U)=0.
\]

\begin{authornote}
\textbf{潜在疑问：}
“将 Maxwell 耦合系统写成 \(\Nop(\bm U)=0\)，如何帮助我们分析微扰？”

\textbf{作者补充：}
具体方程中的时间导数、空间导数和非线性耦合可以合记为一个作用在未知场上的算子。这样无需改变原问题，就能统一讨论解附近的微扰：算子的导数描述微小变化如何传递，也给 Newton 迭代提供每一步要近似求解的线性方程。

\textbf{核心概念联系：}
\[
\text{coupled field equations}
\longrightarrow
\Nop(\bm U)=0
\longrightarrow
\text{linearized operator}
\]
\end{authornote}

在基态 \(\bm U_0\) 附近写
\[
\bm U=\bm U_0+\epsilon\bm u,
\qquad 0<\epsilon\ll1.
\]
若算子足够光滑，则
\[
\Nop(\bm U_0+\epsilon\bm u)
=
\Nop(\bm U_0)
+\epsilon D\Nop[\bm U_0](\bm u)
+\frac{\epsilon^2}{2}
D^2\Nop[\bm U_0](\bm u,\bm u)
+\cdots.
\]
若基态满足
\[
\Nop(\bm U_0)=0,
\]
保留一阶项得
\[
\Lop\bm u=0,
\qquad
\Lop:=D\Nop[\bm U_0].
\]

这里 \(D\Nop[\bm U_0]\) 是 Fr\'echet 导数，也可称线性化算子、Jacobian 算子或切线算子。它是非线性算子在基态处的局部线性近似。

例如非线性介质满足
\[
\bm D=\bm D(\bm E),
\qquad
\bm B=\bm B(\bm H).
\]
令
\[
\bm E=\bm E_0+\epsilon\bm e,
\qquad
\bm H=\bm H_0+\epsilon\bm h.
\]
本构关系的一阶展开为
\[
\bm D(\bm E_0+\epsilon\bm e)
=
\bm D(\bm E_0)
+
\epsilon
\left.
\frac{\partial\bm D}{\partial\bm E}
\right|_{\bm E_0}\bm e
+\Order(\epsilon^2),
\]
\[
\bm B(\bm H_0+\epsilon\bm h)
=
\bm B(\bm H_0)
+
\epsilon
\left.
\frac{\partial\bm B}{\partial\bm H}
\right|_{\bm H_0}\bm h
+\Order(\epsilon^2).
\]

\begin{askai}
AI 展开问题：
\begin{quote}
把有限维 Taylor 展开和 Fr\'echet 导数并排比较，逐步解释为什么非线性材料关系
\[
\bm D=\bm D(\bm E),\qquad \bm B=\bm B(\bm H)
\]
在基态附近会产生切线材料张量。请用一个各向异性材料的矩阵例子说明。
\end{quote}

进一步展开：
\begin{quote}
如果 \(\bm D(\bm E)\) 是非线性的，为什么扰动方程里出现的是 \(\partial\bm D/\partial\bm E|_{\bm E_0}\) 而不是原来的 \(\bm D\)？
\end{quote}
\end{askai}

定义切线材料张量
\[
\bm\varepsilon_{\mathrm{tan}}
=
\left.\frac{\partial\bm D}{\partial\bm E}\right|_{\bm E_0},
\qquad
\bm\mu_{\mathrm{tan}}
=
\left.\frac{\partial\bm B}{\partial\bm H}\right|_{\bm H_0}.
\]
对无色散的瞬时本构，以下进一步取时间不变的基态，使切线材料张量不随时间变化，并取零扰动电流。此时一阶扰动满足线性 Maxwell 型系统
\[
\nabla\times\bm e
=
-\bm\mu_{\mathrm{tan}}\frac{\partial\bm h}{\partial t},
\]
\[
\nabla\times\bm h
=
\bm\varepsilon_{\mathrm{tan}}\frac{\partial\bm e}{\partial t}.
\]

\section{Vlasov 微扰与波的色散关系}
在 Vlasov--Maxwell 系统中，取满足零阶方程的平衡 $f_{s0}$、$\bm E_0=0$ 与给定 $\bm B_0$，设
\[
f_s=f_{s0}+\epsilon f_{s1},
\qquad
\bm E=\epsilon\bm E_1,
\qquad
\bm B=\bm B_0+\epsilon\bm B_1.
\]
忽略 \(\Order(\epsilon^2)\) 项，可得
\[
\frac{\partial f_{s1}}{\partial t}
+\bm v\cdot\nabla_{\bm x}f_{s1}
+\frac{q_s}{m_s}
(\bm v\times\bm B_0)\cdot\nabla_{\bm v}f_{s1}
=
-\frac{q_s}{m_s}
(\bm E_1+\bm v\times\bm B_1)
\cdot\nabla_{\bm v}f_{s0}.
\]
扰动电流为
\[
\bm J_1=\sum_s q_s\int \bm v f_{s1}\dd\bm v.
\]

对常系数、有限个振幅未知量构成的封闭线性系统，或已经完成速度积分响应的动理学系统，进一步采用平面波
\[
\bm u_1(\bm x,t)
=
\widehat{\bm u}
\exp[\ii(\bm k\cdot\bm x-\omega t)],
\]

\begin{authornote}
\textbf{潜在疑问：}
“为什么分析扰动时应该想到平面波？这是不是额外假设了所有解都是波？”

\textbf{作者补充：}
均匀背景附近的常系数线性化算子保持各 Fourier 模态的波数，因此可分别考察每个模态。缓慢变化的背景可在局部冻结系数，以平面波探测局部响应；边界和全局非均匀性则通过相应的整体问题处理。

\textbf{核心概念联系：}
\[
\text{linear perturbation equation}
\longrightarrow
\text{plane-wave ansatz}
\longrightarrow
\text{dispersion relation}
\]
\end{authornote}

则
\[
\partial_t\mapsto-\ii\omega,
\qquad
\nabla\mapsto\ii\bm k.
\]
线性偏微分方程组变为
\[
\bm M(\omega,\bm k)\widehat{\bm u}=0.
\]

\begin{askai}
AI 展开问题：
\begin{quote}
从平面波代换
\[
\partial_t\mapsto-\ii\omega,\qquad \nabla\mapsto\ii\bm k
\]
开始，选一个简单线性 PDE 系统，逐步演示它如何变成
\[
\bm M(\omega,\bm k)\widehat{\bm u}=0.
\]
请解释为什么非零振幅要求 \(\det\bm M(\omega,\bm k)=0\)。
\end{quote}

进一步展开：
\begin{quote}
用一维声波方程或真空 Maxwell 平面波做例子，完整求出对应的色散关系。
\end{quote}
\end{askai}

存在非零解的条件
\[
\det\bm M(\omega,\bm k)=0
\]
就是色散关系。

对仍保留连续速度变量的动理学方程，首先得到的是速度空间响应问题；需指定解析延拓、积分路径与场方程后，才能建立相应介电响应及色散条件。

微扰法按小参数逐阶构造解，稳定性分析考察线性扰动增长率。Fr\'echet 导数也生成非线性代数迭代所用的 Jacobian，其求解算法在第 0 卷讨论。

\chapter{尺度分析与渐近约化}\label{ch:4}

\begin{chaptersummary}
已有的动理学、连续介质和电磁模型包含不同尺度与机制。如何在具体问题中取舍这些机制，是本章的主线。无量纲化将相对量级写成尺度比，渐近分析再据此构造极限模型与边界层。这使前面的模型族获得明确的适用区域，也为后续数值误差与分辨率选择提供依据。
\end{chaptersummary}
\learninggoals{由特征尺度形成无量纲参数；用主导平衡确定内外区域；说明准中性、强磁场和低频约化所去除的快尺度。}
\studyroute{这里略去匹配条件的构造和误差一致性证明。}{singular perturbation; dominant balance; matched asymptotic expansions}{对 $\epsilon u''+u'=0$、$u(0)=0,u(1)=1$ 构造内外解，核对边界层位置并给出一致近似。}

\section{无量纲化、尺度排序与渐近展开}\index{奇异摄动}
常见无量纲数包括
\[
\mathrm{Kn}=\frac{\lambda}{L},
\qquad
\mathrm{Re}=\frac{\rho UL}{\mu},
\qquad
\mathrm{Ma}=\frac{U}{c}.
\]
这些参数度量不同物理机制的相对量级。模型选择本质上是决定哪些项在给定极限下保留。

\begin{authornote}
\textbf{潜在疑问：}
“为什么连续介质方程之后要立刻谈无量纲数？这些数和模型选择有什么因果关系？”

\textbf{作者补充：}
不同物理量的数值大小依赖单位，不能直接比较。选定问题的特征长度、速度和时间后，无量纲化把方程各项的相对大小集中到无量纲参数中。参数的大小提示哪些机制可能主导，以及在哪些极限下可以建立简化模型；尺度的选取仍需结合具体问题。

\textbf{核心概念联系：}
\[
\text{continuum equations}
\longrightarrow
\text{dimensionless ratios}
\longrightarrow
\text{model regimes}
\]
\end{authornote}

Taylor 展开从局部导数描述函数变化，渐近展开则组织小参数极限下的尺度层次。先看 Taylor 形式：
\[
f(x+h)
=
f(x)+hf'(x)+\frac{h^2}{2}f''(x)+\cdots.
\]

\begin{askai}
AI 展开问题：
\begin{quote}
请比较 Taylor expansion 和 asymptotic expansion。给出一个渐近展开有用但级数不收敛的例子，并解释为什么“每次截断后余项更小”比“无限和收敛”更适合某些扰动问题。
\end{quote}

进一步展开：
\begin{quote}
用一个简单函数展示最优截断：计算不同阶截断误差，并说明为什么项最小处附近应该停止。
\end{quote}
\end{askai}

它特别适合分析离散算子与连续算子的局部差异。

渐近展开更一般。设 \(\epsilon\to0\)，若尺度函数满足
\[
\phi_{n+1}(\epsilon)=o(\phi_n(\epsilon)),
\]
且对每个固定 \(N\) 都有
\[
f(\epsilon)
-
\sum_{n=0}^{N}a_n\phi_n(\epsilon)
=
o(\phi_N(\epsilon)),
\]
则写作
\[
f(\epsilon)\sim\sum_{n=0}^{\infty}a_n\phi_n(\epsilon).
\]

渐近展开要求：对每个固定截断阶数，余项比最后保留的尺度更小。级数本身可以发散，尺度也可以包含分数幂、对数与指数项，例如
\[
1,\quad
\epsilon^{1/2},\quad
\epsilon\log\epsilon,\quad
\ee^{-1/\epsilon}.
\]

因此，Taylor 级数是由局部导数产生、以普通幂为基的一类特殊展开；渐近展开强调的是不同项的尺度排序。

若系数增长很快，渐近级数可能先改善后恶化。固定 \(\epsilon\) 后，项
\[
|a_n\epsilon^n|
\]
可能先减小再增大。此时应在最小项附近截断，而不是追求无限求和。这称为最优截断。

若
\[
u=u_0+\epsilon u_1+\epsilon^2u_2+\cdots
\]
代入后，各阶问题保持原方程的微分阶数和边界条件结构，称为正则摄动。

若小参数乘在最高阶导数前，例如
\[
\epsilon u''+u'=0,
\]
则 \(\epsilon\to0\) 时方程阶数下降，部分边界条件无法由外部解满足。这是奇异摄动，通常需要边界层、匹配渐近展开、多尺度法或 WKB 方法。

\section{等离子体示例：小 Debye 长度与准中性极限}

考虑一维无量纲化的 Poisson 方程
\[
-\epsilon^2\frac{\partial^2\phi}{\partial x^2}
=
n_i-n_e,
\qquad
\epsilon=\frac{\lambda_D}{L}\ll1.
\]
这里 \(\lambda_D\) 是 Debye 长度，\(L\) 是宏观尺度。

若直接令 \(\epsilon=0\)，得到
\[
n_i=n_e,
\]
即准中性条件。这说明在主体区域，净电荷密度是高阶小量。

但如果边界给定的电势与主体准中性解不相容，最高阶导数前的小参数使问题成为奇异摄动。主体区域的外部展开可写成
\[
\phi^{\mathrm{out}}
=
\phi_0+\epsilon\phi_1+\epsilon^2\phi_2+\cdots,
\]
对电势变化与归一化净电荷均为 $O(1)$ 的层，平衡 $\epsilon^2/\delta^2=O(1)$ 得到 $\delta=O(\epsilon)$。在这一排序下，边界附近引入拉伸坐标
\[
\xi=\frac{x}{\epsilon}.
\]

\begin{authornote}
\textbf{潜在疑问：}
“为什么边界附近必须引入 \(\xi=x/\epsilon\)，而不是换成任意新的坐标？”

\textbf{作者补充：}
当小参数使最高阶导数项在主体区域看似可以忽略时，降阶后的方程可能无法满足全部边界条件。边界附近的解必须在较短距离内迅速变化，使原本系数很小的导数项重新与其他项达到同一量级。先以未知厚度估计各项并作主导平衡，再据此引入拉伸坐标，才能解析这一薄层。

\textbf{核心概念联系：}
\[
\text{singular perturbation}
\longrightarrow
\text{boundary-layer scale}
\longrightarrow
\xi=x/\epsilon
\]
\end{authornote}

在该尺度下
\[
\frac{\partial}{\partial x}
=
\frac{1}{\epsilon}\frac{\partial}{\partial\xi},
\]

\begin{askai}
AI 展开问题：
\begin{quote}
先以 \(\delta(\epsilon)\) 表示未知的 boundary-layer width，通过 dominant balance 推导 \(\delta\) 随 \(\epsilon\) 的缩放。再比较 stretched coordinate \(x/\epsilon\) 与这一尺度，解释为何令 \(\xi=x/\epsilon\)。
\end{quote}

进一步展开：
\begin{quote}
把外部解、内部解和匹配区域分别画成一张概念图，并说明准中性条件在哪个区域有效。
\end{quote}
\end{askai}

Poisson 方程恢复为
\[
-\frac{\partial^2\phi^{\mathrm{in}}}{\partial\xi^2}
=
n_i-n_e.
\]

在上述排序下，Debye 尺度的电荷分离层承担壁面电势与主体解之间的匹配。求其形状还需粒子响应、通量与壁面条件；其他电势幅度或输运尺度可能产生额外的过渡层。

该例展示了渐近展开的几个典型特征：
\begin{enumerate}
  \item 小参数来自明确的尺度比 \(\lambda_D/L\)；
  \item 主体区域和边界层使用不同坐标；
  \item 令小参数为零会降低方程阶数；
  \item 内外解必须在重叠区域匹配；
  \item 极限模型的有效性依赖于空间区域和边界条件。
\end{enumerate}

在等离子体文献中，类似结构还出现在低频极限、强磁场极限、漂移动理学、回旋动理学和 MHD 极限中。阅读时应先问：小参数是什么，哪些变量被假设为同阶，极限是否降低方程阶数，是否存在快时间、短空间或边界层尺度。

常见记号应按以下方式理解：
\[
f=g+\Order(\epsilon^p)
\]
表示余项被常数倍 \(\epsilon^p\) 控制；
\[
f=g+o(\epsilon^p)
\]
表示余项与 \(\epsilon^p\) 的比值趋于零；
\[
f\sim g
\]
既可能表示 \(f/g\to1\)，也可能表示逐阶渐近展开，必须结合上下文判断。

\section{强磁场、低频与回旋尺度}
带电粒子的回旋频率与典型回旋半径为
\[
\Omega_s=\frac{|q_s|B_0}{m_s},\qquad
\rho_{Ls}=\frac{v_{\mathrm{th},s}}{\Omega_s}.
\]
当 $\rho_{Ls}/L\ll1$ 且 $\omega/\Omega_s\ll1$，可按快回旋角与慢导引中心运动组织变量。
漂移动理学进一步取 $k_\perp\rho_{Ls}\ll1$；回旋动理学允许扰动满足 $k_\perp\rho_{Ls}=O(1)$，
保留有限回旋半径效应。第~\ref{ch:19}~章给出五维分布、回旋平均与自洽场方程。

\part{连续方程的数学结构}
\chapter{函数空间、Sobolev 空间与弱解}\label{ch:5}
\begin{chaptersummary}
前四章给出了物理模型及其尺度假设。求解这些方程还需规定导数、边界值和解的含义。本章从 Poisson 分部积分进入弱导数与 Sobolev 空间，用能量估计说明解为何存在且稳定，再以守恒律说明间断解的选择。下一章将在这些空间上构造有限维残差约束。
\end{chaptersummary}
\learninggoals{从强形式得到含边界项的弱形式；定义弱导数与迹；在给定假设下使用 Poincar\'e 不等式和 Lax--Milgram 定理；区分分片空间与整体空间。}
\studyroute{正文给出弱形式、空间定义和核心估计，略去一般迹定理、紧嵌入及低正则性椭圆理论的证明。}{weak derivative; Sobolev trace; Poincar\'e inequality; Lax--Milgram theorem}{对单位区间的 Poisson 问题逐项分部积分，再用能量估计说明 $H^{-1}$ 数据如何控制 $H_0^1$ 解。}

\section{从强形式到积分恒等式}\index{弱形式}
取有界 Lipschitz 区域 $\Omega\subset\R^d$。先考虑足够光滑的解
\[
-\Delta u=f\quad\text{于 }\Omega,\qquad u=0\quad\text{于 }\partial\Omega.
\]
乘测试函数 $v$ 并积分，散度定理给出
\begin{equation}\label{eq:poisson-green}
\int_\Omega\nabla u\cdot\nabla v\dd x
=\int_\Omega fv\dd x+\int_{\partial\Omega}(\partial_nu)v\dd S.
\end{equation}
令测试函数的边界值为零后，边界项消失。该恒等式只使用 $u$ 的一阶导数，
因而可以在比经典二阶可微函数更大的空间中解释。
非齐次 Dirichlet 数据通过给定提升 $u_g$ 写作 $u=u_g+w$，测试空间仍取零迹空间。
Neumann 数据则进入右端边界配对。

\conceptbridge{为何降低导数阶数后，还要专门规定函数空间？}
{分部积分改变了方程对正则性的要求。平方可积的一阶弱导数使能量积分有意义；迹规定边界条件。函数空间同时决定允许的解、测试方向以及稳定性估计的度量。}
{分部积分 $\longrightarrow$ 弱导数与边界迹 $\longrightarrow$ 变分问题。}

\section{弱导数与空间的完备性}\index{Sobolev 空间}
对局部可积的 $u$，若
\[
\int_\Omega u\,\partial_i\varphi\dd x=-\int_\Omega g_i\varphi\dd x
\quad\text{对所有 }\varphi\in C_c^\infty(\Omega),
\]
则 $g_i$ 是 $u$ 的第 $i$ 个弱导数。分布导数允许右端是一般分布；
例如一维阶跃函数的导数是 Dirac 分布，它不属于 $L^2$。

赋范空间中的每个 Cauchy 列都在该空间收敛时，称其为 Banach 空间。
范数由内积诱导且完备时，称其为 Hilbert 空间。完备性保证能量估计所得的极限仍是允许的函数。
\[
L^2(\Omega)=\{u:\int_\Omega|u|^2\dd x<\infty\},\qquad
(u,v)_{L^2}=\int_\Omega uv\dd x.
\]
这里按几乎处处相等识别函数。定义
\[
H^1(\Omega)=\{u\in L^2(\Omega):\partial_i u\in L^2(\Omega),\ 1\le i\le d\},
\qquad
\norm{u}_{H^1}^2=\norm{u}_{L^2}^2+\norm{\nabla u}_{L^2}^2.
\]
整数 $m\ge0$ 时，$H^m$ 要求所有 $|\alpha|\le m$ 的弱导数属于 $L^2$。
实数阶数可在 $\R^d$ 上通过 Fourier 变换定义：
\[
\norm{u}_{H^s(\R^d)}^2
=\int_{\R^d}(1+|\bm k|^2)^s|\widehat u(\bm k)|^2\dd\bm k,
\]
其中采用酉 Fourier 归一化；区域上的 $H^s$ 可由限制和延拓构造。

\section{迹、Poincar\'e 不等式与能量}\index{迹算子}
对上述区域，迹算子 $\gamma:H^1(\Omega)\to H^{1/2}(\partial\Omega)$
连续地延拓光滑函数的边界限制。定义
\[
H_0^1(\Omega)=\overline{C_c^\infty(\Omega)}^{H^1}
=\{u\in H^1(\Omega):\gamma u=0\}.
\]
对 $v\in H_0^1(\Omega)$，Poincar\'e 不等式给出
\[
\norm{v}_{L^2}\le C_P\norm{\nabla v}_{L^2}.
\]
因此 $\norm{\nabla v}_{L^2}$ 在该空间上是与 $H^1$ 范数等价的能量范数。
纯 Neumann 问题保留常数零空间，需固定平均值，并满足
$\int_\Omega f\dd x+\int_{\partial\Omega}g_N\dd S=0$。

\section{弱问题与 Lax--Milgram 定理}
令 $V=H_0^1(\Omega)$，$V'=H^{-1}(\Omega)$。弱 Poisson 问题为
\[
\text{求 }u\in V:\qquad b(u,v)=\ell(v)\quad\forall v\in V,\qquad
b(u,v)=\int_\Omega\nabla u\cdot\nabla v\dd x.
\]
当 $f\in L^2$ 时 $\ell(v)=\int fv$；一般 $f\in H^{-1}$ 时使用对偶配对。
\begin{theorem}[Lax--Milgram]
设 $V$ 为实 Hilbert 空间，双线性形式满足
$|b(w,v)|\le M\norm{w}_V\norm{v}_V$ 及
$b(v,v)\ge\alpha\norm{v}_V^2$，其中 $\alpha>0$。
若 $\ell\in V'$，则存在唯一 $u\in V$ 满足 $b(u,v)=\ell(v)$，并有
$\norm{u}_V\le\norm{\ell}_{V'}/\alpha$。
\end{theorem}
Poisson 问题以能量范数度量时 $M=\alpha=1$。
取 $v=u$ 即得到该稳定性估计；存在性可通过 Riesz 表示与有界强制算子证明。
对流主导、鞍点或波动问题需要与其算子相符的估计。

\aiexercise{在 $(0,1)$ 上取 $u(x)=x(1-x)$。从强形式 $-u''=2$ 推导弱形式，计算能量与右端配对，再取一个有折点的连续分片线性函数，验证其属于 $H_0^1$。}
{以一维阶跃函数为例计算分布导数；说明它为何不属于 $H^1$。随后推导纯 Neumann 问题的兼容条件与常数零空间。}

\section{守恒律的弱解与间断}
对一维 $u_t+\partial_xF(u)=0$，测试函数
$\varphi\in C_c^\infty([0,T)\times\R)$ 给出
\[
\int_0^T\int_\R(u\varphi_t+F(u)\varphi_x)\dd x\dd t
+\int_\R u_0(x)\varphi(0,x)\dd x=0.
\]
移动间断 $x=s t$ 满足 Rankine--Hugoniot 条件
\[
s(u_R-u_L)=F(u_R)-F(u_L).
\]
弱恒等式提供跨间断的守恒关系。对一维标量守恒律，在适当数据条件下，一族凸熵不等式选择唯一熵解；
熵对 $\eta,q$ 满足 $q'=\eta'F'$，要求分布意义的
$\partial_t\eta(u)+\partial_xq(u)\le0$。有限体积通量将在第~\ref{ch:8}~章近似这一传播过程。

\section{分片 Sobolev 空间}
设 $\mathcal T_h$ 为网格，
\[
H^1(\mathcal T_h)=\{v\in L^2(\Omega):v|_K\in H^1(K),\ K\in\mathcal T_h\}.
\]
其导数先逐单元计算，界面迹可有跳跃。分布导数还包含界面贡献；
DG 的数值通量或椭圆罚项负责这些界面作用。
分片空间规定允许的不连续性，算法仍需进一步给出残差与界面耦合。

\chapter{投影、插值与 Galerkin 框架}\label{ch:6}
\begin{chaptersummary}
弱形式规定了解与测试函数所在的连续空间，离散计算则进一步选择有限维表示及残差约束。本章据此比较插值、正交投影、Galerkin、配置点与控制体守恒，再用误差正交性联系空间逼近与数值解。后续各方法将在这一共同框架下定义自由度、算子和界面耦合。
\end{chaptersummary}
\learninggoals{写出试探与测试空间；区分五类逼近或残差机制；说明能量投影的适用条件；由局部基函数构造质量矩阵、残差和稀疏耦合。}
\studyroute{这里给出强制问题的基本误差估计，略去一般 Petrov--Galerkin 的离散 inf--sup 理论和非协调误差估计。}{Galerkin orthogonality; Riesz representation; Cea lemma; Petrov--Galerkin inf-sup}{对 $-\!u''=f$ 比较两种线性有限元网格的能量误差，再以对流问题解释测试空间改变后的稳定性问题。}

\section{从连续算子到残差约束}
设连续问题为 $\mathcal A u=f$，允许的场属于 $V$，测试空间为 $W$。
弱形式可写为 $b(u,v)=\ell(v)$。
选择有限维试探空间 $V_h$ 和测试空间 $W_h$，得到
\begin{equation}\label{eq:galerkin-general}
\text{求 }u_h\in V_h:\qquad b(u_h,v_h)=\ell(v_h)
\quad\forall v_h\in W_h.
\end{equation}
残差泛函 $r_h(v)=\ell(v)-b(u_h,v)$ 在 $W_h$ 上为零；
它在其余测试方向上通常仍有误差。相同空间称为 Galerkin；
不同试探和测试空间构成 Petrov--Galerkin。对非线性问题，残差与解的关系相应成为非线性。

\conceptbridge{有限维表示确定之后，还缺少哪一步才能得到数值方程？}
{基函数只给出了候选解的表示。数值方程还需规定哪些测试方向、网格点或控制体上的残差为零。这一约束机制与界面处理共同定义离散算子。}
{表示空间 $\longrightarrow$ 残差约束 $\longrightarrow$ 有限维方程。}

\section{插值与正交投影}\index{Galerkin 方法}\index{插值}
插值算子 $I_h$ 通过指定自由度 $\lambda_i$ 满足
$\lambda_i(I_hu)=\lambda_i(u)$。节点插值要求点值有意义；
在一般高维 $H^1$ 空间中，点值未必是连续泛函，可改用局部平均型自由度。

Hilbert 空间内的正交投影 $P_h$ 则满足
\[
P_hu\in V_h,\qquad (u-P_hu,v_h)_V=0\quad\forall v_h\in V_h.
\]
对 $u(x)=x^2$、$V_h=\Poly_1(-1,1)$，端点线性插值得到 $I_hu=1$，
而 $L^2$ 正交投影得到 $P_hu=1/3$。
同一个表示空间可以对应不同的逼近算子。

\section{Riesz 表示、能量范数与误差估计}
Hilbert 空间的 Riesz 映射 $R_V:V\to V'$ 由
$\langle R_Vw,v\rangle=(w,v)_V$ 定义，
使每个连续线性泛函对应唯一的内积表示向量。具体 PDE 的双线性形式还需核对对称性和强制性。
对 $V_h\subset V$ 的协调 Galerkin 近似，连续与离散方程相减得
\[
b(u-u_h,v_h)=0\quad\forall v_h\in V_h.
\]
若 $b$ 对称正定，它定义能量范数 $\norm{v}_b^2=b(v,v)$，并有
\[
\norm{u-u_h}_b=\inf_{w_h\in V_h}\norm{u-w_h}_b.
\]
对有界且强制、但未必对称的 $b$，由任意 $w_h\in V_h$ 可得
\[
\alpha\norm{u-u_h}_V^2
\le b(u-u_h,u-w_h)
\le M\norm{u-u_h}_V\norm{u-w_h}_V.
\]
约去误差范数即得到 C\'ea 估计
\[
\norm{u-u_h}_V\le\frac{M}{\alpha}
\inf_{w_h\in V_h}\norm{u-w_h}_V.
\]
这一结论把数值误差联系到空间逼近能力。Petrov--Galerkin 与非强制问题需离散 inf--sup 条件。

\section{不同方法的自由度与残差机制}
\begin{center}\small
\begin{tabularx}{\textwidth}{@{}l>{\raggedright\arraybackslash}X>{\raggedright\arraybackslash}X@{}}
\toprule
方法 & 表示与自由度 & 方程或残差约束\\
\midrule
FDM & 网格函数、节点值 & 差分算子在网格点满足方程\\
FVM & 单元平均与局部重构 & 每个控制体的积分守恒\\
FEM & 局部基函数、迹相容自由度 & 对有限维测试空间满足弱形式\\
谱 Galerkin & 全局模式系数 & 对所选全局测试模式残差为零\\
谱配置 & 全局插值多项式的节点值 & 配置点满足微分方程\\
DG / FR & 单元局部多项式 & 弱残差 / 修正通量导数及共享界面通量\\
LBM & 各离散速度的粒子群密度 & 迁移与碰撞及其守恒矩\\
\bottomrule
\end{tabularx}
\end{center}
这些方法可共享有限维表示、残差和结构分析的语言。
Hilbert 正交投影的具体解释由内积、试探空间与残差条件共同决定。

\aiexercise{在 $[-1,1]$ 上分别求 $x^2$ 到 $\Poly_1$ 的端点插值与 $L^2$ 投影，显式计算两个误差。再说明对称 Poisson Galerkin 的正交性使用哪一个内积。}
{比较中心差分、单元平均 FVM 和线性 FEM 对同一 Poisson 问题的自由度及残差约束；列出均匀网格上产生相同矩阵所需的积分和边界条件。}

\section{有限维代数系统的形成}
偏微分方程离散后通常得到
\[
\bm A\bm x=\bm b
\]
或非线性残差系统
\[
\bm R(\bm U)=0.
\]

\begin{authornote}
\textbf{潜在疑问：}
“为什么一个连续 PDE 离散后会变成有限维向量方程？自由度、残差和矩阵条目分别来自哪里？”

\textbf{作者补充：}
连续场在每一点都有取值，计算机只能保存有限个参数。离散化先用网格值或基函数系数表示近似场，再通过节点配点、控制体守恒或对测试函数积分等方式，要求方程在有限个意义下成立。于是连续问题转化为有限维代数方程；非线性或含时问题还需进一步线性化或时间离散。

\textbf{核心概念联系：}
\[
\text{continuous PDE}
\longrightarrow
\text{finite-dimensional DOFs}
\longrightarrow
\text{algebraic system}
\]
\end{authornote}

局部微分算子通常产生稀疏矩阵，边界积分算子常产生稠密矩阵。矩阵结构直接反映物理耦合范围。

\begin{authornote}
\textbf{潜在疑问：}
“为什么 differential operator 的局部性会变成 matrix sparsity？为什么积分方程又会变稠密？”

\textbf{作者补充：}
矩阵的一行记录一个离散方程依赖哪些自由度。局部微分算子配合局部基函数或有限差分模板时，该方程通常只涉及附近自由度，因此每行只有少量非零项。若采用全局基函数或积分算子，一个自由度可能与许多远处自由度耦合，矩阵就可能变稠密；稀疏性取决于算子与离散方式的共同选择。

\textbf{核心概念联系：}
\[
\text{operator locality}
\longrightarrow
\text{matrix sparsity or density}
\]
\end{authornote}


取基函数 $\phi_j$，$u_h=\sum_jU_j\phi_j$，式~\eqref{eq:galerkin-general}~产生
$A_{ij}=b(\phi_j,\psi_i)$ 与 $b_i=\ell(\psi_i)$。
对含时问题，
\[
M_{ij}=(\phi_j,\psi_i),\qquad M_h\dot U_h=R_h(U_h,t).
\]
质量矩阵表示时间导数的配对，残差包含体积、边界及源项。
它们的结构由表示与约束产生；代数求解算法在第 0 卷展开。

\chapter{有限差分法：一致性、稳定性与波动误差}\label{ch:7}
\begin{chaptersummary}
函数空间与残差框架确定后，本章用网格函数和差分算子构造最直接的离散模型。Taylor 展开给出一致性，误差演化给出稳定性，修正方程与数值频率分别量化耗散和色散。这些分析将为有限体积、高阶方法及时间推进提供比较尺度。
\end{chaptersummary}
\learninggoals{从模板推导截断误差；说明周期 Fourier 分析的前提；区分稳定性、CFL 必要条件和收敛性；计算振幅与相位误差。}
\studyroute{正文通过平流、网格范数和误差传播建立差分分析；有限边界 SBP--SAT 分析与完整收敛证明留待进一步展开。}{grid function; truncation error; von Neumann analysis; modified equation; SBP--SAT}{对上风格式推导逐模态放大因子，再以离散 Parseval 公式证明周期离散 L2 收缩。}
\section{网格函数与差分算子}
均匀节点 $x_i=ih$ 上的自由度为 $U_i\approx u(x_i)$。定义
\[
(D_-U)_i=\frac{U_i-U_{i-1}}h,\qquad
(D_0U)_i=\frac{U_{i+1}-U_{i-1}}{2h},\qquad
\norm{U}_{2,h}^2=h\sum_i|U_i|^2.
\]
网格函数属于有限维序列空间。差分算子在网格点近似微分方程；
插值可用于比较网格函数与连续场。
\section{一致性、稳定性与修正方程}\index{有限差分法}
考虑线性平流方程
\[
u_t+au_x=0,
\qquad
u(x,t)=u_0(x-at).
\]

\begin{authornote}
\textbf{潜在疑问：}
“前面讲的是复杂物理模型，为什么进入数值离散时突然只研究线性平流方程？”

\textbf{作者补充：}
复杂输运方程同时包含多种物理机制，难以辨认误差来自哪里。线性平流方程只保留以固定速度传播这一基本过程，因而可以单独考察离散格式如何改变传播、振幅和相位。它是分析数值机制的参照模型，其结论仍需结合非线性、变系数和边界等因素重新判断。

\textbf{核心概念联系：}
\[
\text{complicated transport physics}
\longrightarrow
\text{linear advection as minimal model}
\]
\end{authornote}

前向时间差分满足
\[
\frac{u_i^{n+1}-u_i^n}{\Delta t}
=
u_t+\Order(\Delta t),
\]
中心空间差分满足
\[
\frac{u_{i+1}-u_{i-1}}{2\Delta x}
=
u_x+\Order(\Delta x^2).
\]
离散格式的阶数来自离散算子与连续算子的局部展开匹配。

对线性常系数格式，可令误差模态为
\[
e_i^n=\widehat e^n\ee^{\ii kx_i},
\]

\begin{authornote}
\textbf{潜在疑问：}
“为什么任意误差可以用 Fourier 模态来分析？为什么研究单个模态就能判断整体误差传播？”

\textbf{作者补充：}
在线性常系数、平移不变的格式中，每个离散 Fourier 模态经过一步更新后仍具有相同波数，只改变其复振幅。因此可逐个研究放大因子，判断各模态是否增长。这种分析最直接适用于无限或周期均匀网格；有限边界、变系数和非线性可能引入额外耦合，需要另行分析。

\textbf{核心概念联系：}
\[
\text{arbitrary numerical error}
\longrightarrow
\text{Fourier modes}
\longrightarrow
\text{amplification factor}
\]
\end{authornote}

\begin{askai}
AI 展开问题：
\begin{quote}
从任意离散误差的 Fourier decomposition 开始，逐步建立 von Neumann stability analysis：说明线性常系数差分格式为何保持各 Fourier mode，并据此分别研究其 amplification factor。
\end{quote}

进一步展开：
\begin{quote}
对一阶上风格式完整代入 \(e_i^n=\widehat e^n e^{\ii kx_i}\)，推导 \(G(k)\)，再画出不同 CFL 数下的 \(|G(k)|\)。
\end{quote}
\end{askai}

得到
\[
\widehat e^{n+1}=G(k)\widehat e^n.
\]
单步放大因子刻画各模态变化。对要求无增长的平流离散，收缩条件写成
\[
|G(k)|\le1.
\]

显式格式还受到 CFL 条件约束。定义
\[
\nu=\frac{a\Delta t}{\Delta x}.
\]
其含义是数值依赖域必须覆盖物理依赖域。

\begin{authornote}
\textbf{潜在疑问：}
“为什么 \(\nu=a\Delta t/\Delta x\) 会和依赖域有关？这个数为什么限制显式格式？”

\textbf{作者补充：}
特征速度决定真实解的物理依赖域，离散模板和时间步长决定数值信息的传播范围。网格细化时，数值依赖域覆盖物理依赖域是正确收敛的必要条件；稳定性还需由误差演化估计检验。CFL 条件正是这两种传播范围之间的联系。

\textbf{核心概念联系：}
\[
\text{amplification factor stability}
\longrightarrow
\text{CFL number}
\longrightarrow
\text{domain of dependence}
\]
\end{authornote}

当 \(a>0\) 时，一阶上风格式为
\[
u_i^{n+1}
=
u_i^n-\nu(u_i^n-u_{i-1}^n)
=
(1-\nu)u_i^n+\nu u_{i-1}^n.
\]
在
\[
0\le\nu\le1
\]
时，新值是旧值的凸组合，因此不会产生新的局部极值。

其修正方程为
\[
u_t+au_x
=
\frac{a\Delta x}{2}(1-\nu)u_{xx}
+\cdots.
\]

\begin{askai}
AI 展开问题：
\begin{quote}
对一阶上风格式做 modified equation analysis。请从 Taylor 展开开始逐项推导
\[
u_t+au_x=\frac{a\Delta x}{2}(1-\nu)u_{xx}+\cdots
\]
并解释右侧扩散项如何对应数值耗散。
\end{quote}

进一步展开：
\begin{quote}
写一个最小 Python 数值实验，比较中心格式和上风格式传播一个 Gaussian pulse 时的耗散差异。
\end{quote}
\end{askai}

右侧等效扩散项解释了上风格式的数值耗散：它抑制高频误差，也抹平真实小尺度。
一般线性稳定性要求在固定物理时间 $T$ 内
$\sup_{n\Delta t\le T}\norm{S_h^n}\le C_T$，常数与网格无关。
对平移不变标量更新，Parseval 关系把这一条件转为模态幂的统一界。
允许 $|G|\le1+C\Delta t$ 的有界增长；一次更新的 $|G|>1$ 需结合细化极限判断。
例如
\[
|G(k)|>1,
\]
表示该模态单步增长；若某个随网格细化持续存在的放大值与 1 有固定间隔，固定物理时间内的模态幂会无界增长。
对适定线性初值问题，在相容的离散设置中，一致性与稳定性共同保证收敛。

\section{数值色散与波动误差}
设连续波解为
\[
u=\widehat u\ee^{\ii(kx-\omega t)}.
\]
对线性平流方程，
\[
\omega=ak.
\]
相速度和群速度均为 \(a\)：
\[
c_p=\frac{\omega}{k}=a,
\qquad
c_g=\frac{\dd\omega}{\dd k}=a.
\]

离散格式通常得到数值频率
\[
\omega_{\mathrm{num}}(k).
\]

\begin{askai}
AI 展开问题：
\begin{quote}
选择一个具体差分格式，对一维线性平流方程推导数值频率 \(\omega_{\mathrm{num}}(k)\)。请分别指出 \(\operatorname{Re}\omega_{\mathrm{num}}\) 和 \(\operatorname{Im}\omega_{\mathrm{num}}\) 如何对应相位误差与耗散误差。
\end{quote}

进一步展开：
\begin{quote}
比较两个格式在同一个波数 \(k\) 下的相速度误差，并估计传播很多波长后相位偏差会累积到多大。
\end{quote}
\end{askai}

若
\[
\operatorname{Re}\omega_{\mathrm{num}}\neq ak,
\]
则出现数值色散；若
\[
\operatorname{Im}\omega_{\mathrm{num}}<0,
\]
则出现数值耗散。

流体激波计算通常需要适量耗散，声学、电磁波和弹性波计算则更关注长距离传播中的相位误差。格式选择必须匹配目标物理。

\chapter{有限体积法：数值通量与高分辨率格式}\label{ch:8}
\begin{chaptersummary}
守恒律的弱形式以控制体通量平衡描述间断。本章从单元平均、Riemann 问题和共同通量构造有限体积法，再利用 Godunov 阶数障碍说明非线性重构的作用。守恒、单调性和熵稳定性分别约束通量与更新，为 DG 和 FR 的界面设计提供基础。
\end{chaptersummary}
\learninggoals{由积分守恒构造单元更新；解释 Riemann 通量；区分 Godunov 定理、TVD 与高阶重构；核对重构的平均值和单位。}
\studyroute{正文依次讨论 Godunov、MUSCL、TVD 与 WENO；系统特征分解、完整光滑度指标和非线性熵收敛证明留待进一步展开。}{finite volume; Riemann problem; Godunov order barrier; MUSCL; WENO}{在三个单元平均值上计算 minmod 重构，并将 Rusanov 通量代入两个相邻控制体，核对内部通量的符号。}
\section{控制体、单元平均与共同通量}\index{有限体积法}
有限体积法从控制体守恒式出发：
\[
\frac{\dd}{\dd t}
\int_{\Omega_i}U\dd V
+
\int_{\partial\Omega_i}
\bm F\cdot\bm n\dd S
=
\int_{\Omega_i}S\dd V.
\]
相邻单元以相反符号使用同一界面通量，因此内部通量在整体求和时抵消，得到局部与全局守恒关系。


定义 $\bar U_K=|K|^{-1}\int_KU$。自由度是单元平均，界面状态由平均值重构。
一致数值法向通量满足 $\widehat F(U,U;n)=F(U)\cdot n$，交换两侧并翻转法向后变号。
对一维标量方程，Riemann 初值为界面两侧常数 $U_L,U_R$；
精确解在界面上的通量给出 Godunov 通量，近似求解器则以可计算的波速近似。
例如
\[
\widehat F(U_L,U_R)=\frac{F(U_L)+F(U_R)}2-\frac{\alpha}2(U_R-U_L),
\qquad
\alpha\ge\max_{[U_L,U_R]}|F'(U)|.
\]
线性平流 $F(U)=aU$、$a>0$ 时精确通量是 $aU_L$。
单元更新还需积分通量在一个时间步内的变化；高阶空间重构与时间积分共同决定全离散精度。

\section{Godunov 定理与非线性重构}
Godunov 定理是守恒律数值计算中的基本限制。对线性标量平流问题，任何满足以下条件的线性格式：

\begin{enumerate}
  \item 线性；
  \item 单调，即不产生新的极值；
  \item 一致；
\end{enumerate}

在一般非整数 Courant 数下，线性单调一致格式的精度至多为一阶。对常速平流，整数 Courant 数的精确网格平移是特殊情形。

直观地说，高阶线性重构需要使用正负交替的权重来提高局部逼近阶数，而单调性要求更新具有凸组合性质。两者一般不能同时成立。

因此，提高单调格式的精度需要改变上述条件。现代高分辨率方法通常采用非线性重构：

\begin{authornote}
\textbf{潜在疑问：}
“为什么 Godunov 定理之后自然要引入非线性重构？非线性到底绕开了哪个限制？”

\textbf{作者补充：}
Godunov 定理指出，在线性单调格式的适当条件下，精度至多为一阶。高分辨率方法因此让重构或限制器依赖局部解：光滑区域保留较高阶近似，在陡峭区域调整斜率以抑制振荡。这使格式成为非线性的，从而不再落在定理所约束的线性格式类别中。

\textbf{核心概念联系：}
\[
\text{Godunov theorem}
\longrightarrow
\text{nonlinear reconstruction}
\longrightarrow
\text{high-resolution schemes}
\]
\end{authornote}

\begin{itemize}
  \item 保留守恒性；
  \item 允许算法是非线性的；
  \item 在光滑区域恢复高阶；
  \item 在间断附近退化为一阶单调格式。
\end{itemize}

以有限体积法为例，
\[
\frac{\dd \bar U_i}{\dd t}
=
-\frac{1}{\Delta x}
\left(
\widehat F_{i+1/2}
-
\widehat F_{i-1/2}
\right).
\]
关键在于构造界面左右状态
\[
U_{i+1/2}^{L},
\qquad
U_{i+1/2}^{R},
\]
并通过 Riemann 求解器或近似 Riemann 求解器计算数值通量。

MUSCL 方法采用分片线性重构，例如
\[
U_{i+1/2}^{L}
=
U_i+\frac12\sigma_i,
\]

\begin{askai}
AI 展开问题：
\begin{quote}
用一维有限体积的三个相邻 cell average 构造 MUSCL reconstruction。分别用 unlimited slope 和 minmod limiter 计算 \(U_{i+1/2}^L\)，说明限制器如何在光滑区保留二阶、在间断附近退回一阶。
\end{quote}

进一步展开：
\begin{quote}
给一个含间断的数据例子，逐项计算 minmod、van Leer 和 Superbee 的斜率差异。
\end{quote}
\end{askai}

这里 $\sigma_i$ 表示一个单元宽度上的增量，量纲与 $U_i$ 相同。
例如 $\sigma_i=\operatorname{minmod}(U_i-U_{i-1},U_{i+1}-U_i)$，
其中同号时 minmod 取绝对值较小者，异号时取零。
斜率 \(\sigma_i\) 随局部光滑性调整；minmod、van Leer 和 Superbee 采用不同的限制方式。这样可以在光滑区保留近似二阶精度，在激波或间断附近降低重构阶数以抑制振荡。

TVD 方法要求离散总变差不增加：
\[
TV(U^{n+1})\le TV(U^n),
\qquad
TV(U)=\sum_i|U_{i+1}-U_i|.
\]
它把“不过度产生振荡”转化为可分析的离散条件。

ENO/WENO 方法不直接限制单一斜率，而是在多个候选模板之间选择或加权。WENO 在光滑区域使用高阶组合，在包含间断的模板上自动降低权重。

Godunov 定理揭示的是一个结构性折中：
\[
\text{高阶精度}
\quad\leftrightarrow\quad
\text{单调性与无振荡}.
\]
格式设计据局部解的光滑性调整重构与耗散。

\section{守恒、熵与高阶光滑区}
WENO 以候选重构 $q_r$ 的非线性权重形成界面值
\[
U_{i+1/2}^{L}=\sum_r\omega_rq_r,\qquad
\omega_r=\frac{\alpha_r}{\sum_s\alpha_s},\qquad
\alpha_r=\frac{d_r}{(\epsilon_W+\beta_r)^2}.
\]
$\beta_r$ 度量模板光滑性，$d_r$ 是对应高阶线性组合的理想权重；
$\epsilon_W$ 的尺度需与变量及网格相容。
光滑区阶数依赖权重恢复到理想值的速度，极值附近的行为需单独分析。
标量 TVD 的全离散结论依赖通量、限制器和时间步长；
一般系统的熵稳定设计还需熵变量和相容数值熵通量。WENO 原始构造见~\cite{jiang1996}。

\chapter{有限元、谱方法与积分算子离散}\label{ch:9}
\begin{chaptersummary}
Galerkin 框架将误差分析联系到空间选择与测试条件。本章先装配局部有限元，再比较 Fourier、Legendre 和 Chebyshev 的全局表示，最后以 MoM 构造积分算子的离散。局部或全局支撑决定耦合范围，正则性与求积共同影响误差，因而将逼近理论自然引向第 0 卷的矩阵结构与求解。
\end{chaptersummary}
\learninggoals{装配一维线性有限元质量与刚度矩阵；区分谱 Galerkin、谱配置和伪谱；用积分核构造 MoM 矩阵；解释光滑性对逼近速度的影响。}
\studyroute{略去高维单元族、谱误差定理的一般证明、奇异积分与辐射条件的完整处理；矩阵求解与快速作用算法见第 0 卷。}{Galerkin finite element; spectral collocation; Fourier pseudospectral; Chebyshev approximation; method of moments}{对一维 Poisson 问题分别写出局部线性有限元和全局 Legendre Galerkin 的自由度、边界约束及矩阵条目。}
\section{局部有限元空间与装配}\index{有限元法}
有限元法从弱形式出发。以 Poisson 方程
\[
-\nabla^2u=f
\]
为例，乘测试函数并分部积分：
\[
\int_\Omega\nabla v\cdot\nabla u\dd V
=
\int_\Omega vf\dd V+\text{边界项}.
\]

\begin{askai}
AI 展开问题：
\begin{quote}
对一维 Poisson 方程做二阶中心差分和线性有限元离散，各写出 5 个内部网格点对应的矩阵，并比较 differential operator 的 locality 如何变成 matrix sparsity。
\end{quote}

进一步展开：
\begin{quote}
用同一个 Poisson 例子解释 strong form、weak form、test function、basis function 和 stiffness matrix 之间的关系。
\end{quote}
\end{askai}

展开
\[
u_h=\sum_jU_j\phi_j
\]
后得到
\[
\bm K\bm U=\bm F.
\]

上述展开将弱形式转为基函数系数的方程。下面以一维线性单元显式计算局部矩阵，再由共享节点完成装配。


在一维单元 $[x_j,x_{j+1}]$ 上，长度为 $h_j$，线性形函数分别为
$\phi_0=1-\xi$、$\phi_1=\xi$，$\xi=(x-x_j)/h_j$。
直接积分得到
\[
K^{(j)}=\frac1{h_j}\begin{pmatrix}1&-1\\-1&1\end{pmatrix},
\qquad
M^{(j)}=\frac{h_j}6\begin{pmatrix}2&1\\1&2\end{pmatrix}.
\]
共享节点对应同一自由度，局部条目按节点映射相加。
协调有限元空间满足所需迹条件；$H^1$ 标量元用于 Poisson，
Maxwell 旋度问题则需要相应的切向迹结构，见第二卷。

\conceptbridge{同一组全局多项式为何可以产生 Galerkin 和配置两种谱方法？}
{模式系数与节点值是表示选择。Galerkin 通过积分配对约束残差，配置法在选定节点约束微分残差。两者在适当求积下可能产生相关矩阵，具体关系由测试条件、积分和边界处理决定。}
{全局近似空间 $\longrightarrow$ 模式或节点自由度 $\longrightarrow$ 残差机制。}

\section{全局模式、配置点与伪谱}\index{谱方法}
局部有限元利用紧支撑组织装配，全局谱表示则利用整体光滑性提高精度。周期 $[0,2\pi)$ 上的 Fourier Galerkin 近似取
\[
u_N(x)=\sum_{k=-N}^{N}\widehat u_k\ee^{\ii kx},\qquad
\widehat{(\partial_xu_N)}_k=\ii k\widehat u_k.
\]
Galerkin 以保留的 Fourier 模式测试残差。伪谱计算把模式转到节点，
在物理空间计算非线性乘积，再变回模式；未解析的高波数可能混叠到低波数，
二次非线性可用适当填充或截断消混叠。

区间 $[-1,1]$ 上，Legendre 多项式满足
\[
\int_{-1}^1P_jP_k\dd x=\frac{2}{2k+1}\delta_{jk}.
\]
Chebyshev 多项式 $T_k(\cos\theta)=\cos(k\theta)$
在权 $(1-x^2)^{-1/2}$ 下正交；该权内积与物理无权 $L^2$ 内积需要分别指定。
谱配置可用 Chebyshev--Lobatto 节点 $x_j=\cos(j\pi/p)$，
令 $D_{ij}=\ell_j'(x_i)$，直接在内部节点约束微分残差，边界节点施加边值。

对于足够光滑且解析延拓合适的解，增加全局次数可带来快速甚至指数型收敛。
有限正则性或间断通常只给代数收敛，界面附近还可能出现 Gibbs 振荡。
局部多项式便于几何适配与局部加密，全局模式则利用整体光滑性。
误差比较需报告总自由度、积分误差和时间误差。

\section{积分算子与 MoM}\index{矩量法}
计算电磁学中的矩量法通常从积分方程
\[
\Lop J=g
\]
出发。展开
\[
J=\sum_{n=1}^N I_nf_n,
\]

\begin{authornote}
\textbf{潜在疑问：}
“为什么计算电磁学这里突然从微分方程转向积分方程和电流展开？”

\textbf{作者补充：}
对某些线性散射或辐射问题，可以利用已知背景介质的 Green 函数，把区域内的场表示为边界上等效源或电流产生的场。这样未知量从体域场转到边界密度，并通过边界条件确定；维数可能降低，但 Green 函数的远程作用也使离散矩阵通常变为稠密。

\textbf{核心概念联系：}
\[
\text{field PDE}
\longrightarrow
\text{integral equation}
\longrightarrow
\text{current expansion}
\]
\end{authornote}

并对测试函数 \(w_m\) 约束积分方程残差：
\[
\avg{w_m,\Lop J}
=
\avg{w_m,g}.
\]

\begin{askai}
AI 展开问题：
\begin{quote}
用三段分片常数基函数离散一个简单一维积分方程
\[
\Lop J=g
\]
显式写出矩阵条目 \(\avg{w_m,\Lop f_n}\)，并解释为什么每个测试函数通常会和所有基函数耦合。
\end{quote}

进一步展开：
\begin{quote}
比较微分方程离散矩阵和积分方程离散矩阵的非零结构，并说明快速多极子或层次矩阵想减少的是什么成本。
\end{quote}
\end{askai}

得到
\[
\bm Z\bm I=\bm V.
\]

由于积分核具有长程作用，$\bm Z$ 通常稠密。该结构及快速矩阵作用的算法引向第 0 卷。

具体地，若 $(\mathcal LJ)(x)=\int_\Gamma G(x,y)J(y)\dd S_y$，
则
\[
Z_{mn}=\int_\Gamma w_m(x)\int_\Gamma G(x,y)f_n(y)\dd S_y\dd S_x,\qquad
V_m=\int_\Gamma w_mg\dd S.
\]
这里的测试配对属于积分方程离散；其稳定性取决于算子、迹空间、核与基函数。

\part{现代高阶离散算法专题}
\chapter{不连续 Galerkin 方法}\label{ch:10}
\begin{chaptersummary}
有限体积通过共享通量保证单元间守恒，有限元通过测试函数约束残差。DG 将这两个机制结合为局部高阶多项式、逐单元弱形式与界面通量。本章以一维平流推导质量矩阵和局部演化，再由常数与解本身作测试，分别得到守恒与能量恒等式，并进一步讨论扩散的界面罚项。
\end{chaptersummary}
\learninggoals{构造分片试探空间与迹；推导一维 DG 矩阵；证明局部与整体守恒；区分双曲上风稳定与椭圆罚项强制性。}
\studyroute{略去一般非线性熵稳定、多维边界估计和 DG 椭圆误差定理；正文使用常系数、仿射单元及精确积分。}{discontinuous Galerkin; broken polynomial space; upwind flux; symmetric interior penalty; modal nodal DG}{对二次 Legendre 基完整装配平流单元的质量矩阵、体积矩阵和端点向量，并逐界面求和验证能量估计。}

\section{局部空间、自由度与迹}\index{DG 方法}
在周期区间上，单元 $K_j=(x_{j-1/2},x_{j+1/2})$ 的长度为 $h_j$。
取
\[
V_h^p=\{v\in L^2:v|_{K_j}\in\Poly_p(K_j)\}.
\]
每个单元有 $p+1$ 个标量自由度；相邻迹独立。考虑
\[
u_t+a u_x=0,\qquad F(u)=au,\qquad a\in\R.
\]
固定界面正向为从左到右，令 $u^-$、$u^+$ 为两侧迹，
$[u]=u^+-u^-$，$\{u\}=(u^-+u^+)/2$。一致上风通量为
\[
\widehat f=a\{u\}-\frac{|a|}{2}[u].
\]
$a>0$ 时 $\widehat f=au^-$。

\conceptbridge{逐单元允许跳跃后，信息通过什么机制跨单元传播？}
{分片空间允许各单元独立近似；分部积分产生界面通量。用同一个 Riemann 通量连接两侧迹，便将局部残差耦合起来。常数测试函数给出单元守恒，高阶测试函数决定内部多项式的演化。}
{分片多项式 $\longrightarrow$ 逐单元弱形式 $\longrightarrow$ 共享数值通量。}

\section{一维平流的完整弱形式}
取任意 $v_h\in\Poly_p(K_j)$，乘方程并分部积分：
\[
\int_{K_j}v_h\partial_tu_h\dd x
-\int_{K_j}a u_h\partial_xv_h\dd x
+\widehat f_{j+1/2}v_h(x_{j+1/2}^{-})
-\widehat f_{j-1/2}v_h(x_{j-1/2}^{+})=0.
\]
映射 $x=x_j+J_j\xi$，$\xi\in[-1,1]$，$J_j=h_j/2$。
展开 $u_h=\sum_{\ell=0}^pU_{j\ell}\phi_\ell(\xi)$，定义
\[
M_{i\ell}=\int_{-1}^1\phi_i\phi_\ell\dd\xi,\quad
S_{i\ell}=\int_{-1}^1\phi_i'\phi_\ell\dd\xi,\quad
r_i=\phi_i(1),\quad l_i=\phi_i(-1).
\]
对每个基函数测试，得到
\begin{equation}\label{eq:dg-local}
J_jM\dot U_j=a S U_j-r\widehat f_{j+1/2}+l\widehat f_{j-1/2}
=:R_j(U).
\end{equation}
体积矩阵仅作用于本单元，界面残差调用邻居迹。
在 Legendre 模态基中，
$M_{ii}=2/(2i+1)$。对 $p=1$，
\[
M=\begin{pmatrix}2&0\\0&2/3\end{pmatrix},\quad
S=\begin{pmatrix}0&0\\2&0\end{pmatrix},\quad
r=\begin{pmatrix}1\\1\end{pmatrix},\quad
l=\begin{pmatrix}1\\-1\end{pmatrix}.
\]
取 $a>0$，令 $Q_j=U_{j0}+U_{j1}$，则
\[
h_j\dot U_{j0}=-aQ_j+aQ_{j-1},\qquad
\frac{h_j}{3}\dot U_{j1}=2aU_{j0}-aQ_j-aQ_{j-1}.
\]
这给出从自由度到局部 ODE 的可直接实现形式。

\section{局部守恒与能量恒等式}
令 $v_h=1$，
\[
\frac{\dd}{\dd t}\int_{K_j}u_h\dd x
=-\widehat f_{j+1/2}+\widehat f_{j-1/2}.
\]
所有单元求和时，共享通量抵消。$p=0$ 的上风 DG 恢复一阶有限体积空间离散。
令 $v_h=u_h$，对每个界面合并两侧贡献：
\[
\frac12\frac{\dd}{\dd t}\sum_j\int_{K_j}u_h^2\dd x
=\sum_f\left(\widehat f-a\{u_h\}\right)[u_h]
=-\frac{|a|}{2}\sum_f[u_h]^2.
\]
周期边界、常系数与精确积分使等式闭合。中心通量给出零右端。
全离散稳定性还需时间积分器与空间谱相容；有物理边界时能量式需计入边界通量。

\section{模态、节点与椭圆问题}
模态 DG 保存基函数系数；节点 DG 保存插值节点值。
当两者表示同一个多项式空间，且积分与通量相同，基变换连接其半离散方程。
质量集块或求积改变可能产生不同离散算子。

对周期一维 Poisson 问题，可在均值为零的分片空间中使用对称内罚形式
\begin{equation}\label{eq:sip}
\begin{aligned}
b_h(u,v)={}&\sum_K\int_Ku'v'\dd x
+\sum_f\bigl(\{u'\}[v]+\{v'\}[u]\bigr)\\
&+\sum_f\frac{\sigma(p+1)^2}{h_f}[u][v].
\end{aligned}
\end{equation}
此处延用右迹减左迹的标量跳跃，因此一致性项取正号。
$h_f$ 取相邻尺度的适当最小值，$\sigma$ 足够大时，
迹逆估计可将界面交叉项吸收到体积与罚项中，得到离散强制性。
Dirichlet 边界还需相容的边界一致性项与罚项，进一步构造见~\cite{arnold2002}。

\aiexercise{对周期均匀网格和 $p=2$，从式~\eqref{eq:dg-local}~逐项计算 $M,S,l,r$，给出全局块耦合。取常数测试函数及解本身，分别验证总质量和上风能量恒等式。}
{对式~\eqref{eq:sip}~代入光滑精确解并检查界面项符号；再用离散迹不等式说明罚参数需要控制哪些交叉项。}

\chapter{通量重构方法}\label{ch:11}
\begin{chaptersummary}
DG 的强形式可以把界面残差表示为单元内部的提升。FR 从这个思路出发，先插值局部通量，再用端点修正函数接入共同通量。本章推导其节点更新和守恒性质，并通过质量矩阵提升条件给出与特定 DG 离散的精确关系。
\end{chaptersummary}
\learninggoals{由一维守恒律构造不连续通量与修正通量；写出 FR 空间算子；证明单元平均守恒；核对 DG 关系所需的基、积分和通量条件。}
\studyroute{略去多维曲网格、非线性别名误差及一般 FR 熵分析；下一章进一步构造能量稳定修正函数。}{Huynh flux reconstruction; correction function; lifting operator; SBP correction procedure}{在三个 Gauss 节点上构造二次解多项式和三次修正函数，组装节点 FR 算子，并与精确积分的 DG 强形式比较。}

\section{不连续通量与端点修正}\index{FR 方法}
考虑一维 $u_t+\partial_xF(u)=0$。在仿射单元 $x=x_j+J_j\xi$ 上，
选择 $p+1$ 个不同节点 $\xi_i$ 与 Lagrange 基 $\ell_i$，
\[
u_j^\delta(\xi)=\sum_{i=0}^pu_{ji}\ell_i(\xi),\qquad
f_j^D(\xi)=\sum_{i=0}^pF(u_{ji})\ell_i(\xi).
\]
$f^D$ 是通量的次数 $p$ 插值；非线性 $F(u^\delta)$ 的次数可能更高。
界面求解器由两侧解迹产生共同的正向通量 $f_L^*,f_R^*$。
取次数至多 $p+1$ 的修正函数，满足
\[
g_L(-1)=1,\quad g_L(1)=0,\qquad
g_R(-1)=0,\quad g_R(1)=1.
\]
修正通量为
\begin{equation}\label{eq:fr-flux}
f^C=f^D+(f_L^*-f^D(-1))g_L+(f_R^*-f^D(1))g_R.
\end{equation}
于是 $f^C(-1)=f_L^*$、$f^C(1)=f_R^*$。
全局共享的界面通量在每侧局部多项式中得到相同端点值。FR 原始构造见~\cite{huynh2007}。

\conceptbridge{界面通量只给出两个端点值，为何足以改变全部内部自由度？}
{修正函数将端点的通量差延伸为单元内多项式。取其导数后，每个解节点都接收到界面残差。修正函数的形状决定这一提升如何作用于各多项式模式，并影响谱与稳定性。}
{端点通量差 $\longrightarrow$ 修正多项式 $\longrightarrow$ 内部残差提升。}

\section{FR 节点更新与局部守恒}
令 $D_{ik}=\ell_k'(\xi_i)$、$C_{iL}=g_L'(\xi_i)$、$C_{iR}=g_R'(\xi_i)$。
在节点上约束修正后的微分残差，得到
\begin{equation}\label{eq:fr-operator}
J_j\dot u_{ji}
=-\sum_kD_{ik}F(u_{jk})
-(f_L^*-f_L^D)g_L'(\xi_i)
-(f_R^*-f_R^D)g_R'(\xi_i).
\end{equation}
记 $R$ 为取两端点值的 $2\times(p+1)$ 矩阵，则
$J_j\dot{\bm u}_j=-D\bm f_j-C(\bm f^*-R\bm f_j)$。
导数均为次数至多 $p$ 的多项式，故该节点等式在单元内也是多项式恒等式。
精确积分给出
\[
\frac{\dd}{\dd t}\left(J_j\int_{-1}^1u_j^\delta\dd\xi\right)
=-f_R^*+f_L^*.
\]
共同通量与端点条件因此保证单元平均守恒；具体修正函数仍需满足所需稳定性条件。

\section{与 DG 的矩阵关系}
在同一基中，令 $M_{ik}=\int\ell_i\ell_k$，$B=\diag(-1,1)$。
多项式分部积分给出
\[
MD+D^{\mathsf T}M=R^{\mathsf T}BR.
\]
对线性通量 $F(u)=au$，DG 强形式为
\[
J_j\dot{\bm u}_j=-aD\bm u_j-M^{-1}R^{\mathsf T}B
(\bm f^*-aR\bm u_j).
\]
因此，当
\begin{equation}\label{eq:fr-dg-lift}
C=M^{-1}R^{\mathsf T}B,
\end{equation}
且试探表示、积分、映射和通量相同时，FR 与该 DG 半离散算子一致。
模态实现通过可逆基变换给出同一关系。
对非线性通量、欠积分或其他修正函数，需重新比较各个体积与边界算子。

\section{误差与谱分析}
修正函数改变单元模式的界面反馈。周期均匀网格可对 Bloch 符号求谱，
比较物理分支的耗散、色散和显式稳定步长。
能量法则直接检查一个范数下的演化，并纳入物理边界。
耗散、色散与稳定步长还取决于通量、修正函数和时间积分器。按总自由度与误差比较成本时，应将这些选择与形式阶数一并记录。

\aiexercise{对 $p=1$、线性通量与 Gauss 解节点，计算 $M,D,R,B$。由式~\eqref{eq:fr-dg-lift}~求 $g_L',g_R'$，积分并核对四个端点条件，最后与 DG 局部矩阵逐项比较。}
{将一个修正函数改变为另一满足端点条件的多项式。检验单元守恒是否保持，再计算周期符号特征值，说明守恒与稳定性分别需要什么条件。}

\chapter{VCJH 修正函数与能量稳定性}\label{ch:12}
\begin{chaptersummary}
FR 的端点条件给出守恒，但内部修正仍有自由度。本章在一维线性平流中用修正能量范数约束这一自由度，构造 VCJH 单参数族。归一化、质量矩阵与积分规则固定后，可以准确说明其 DG、特定 SD 和 Huynh 修正函数之间的关系。
\end{chaptersummary}
\learninggoals{写出 VCJH 参数化修正函数；推导正定条件与线性能量恒等式；在指定条件下识别 DG 和 SD 参数；以谱和误差评估参数。}
\studyroute{正文证明限于一维常系数平流、仿射静止网格及精确多项式积分；非线性熵、曲网格和扩散扩展作为进一步阅读。}{VCJH correction function; energy stable flux reconstruction; modified Sobolev norm; SBP CPR}{对 $p=1,2,3$ 计算修正函数端点、提升矩阵和修正质量矩阵，并检验线性能量矩阵恒等式。}

\section{参数化修正函数}\index{VCJH 方法}
令 $p\ge1$，$P_p$ 为标准 Legendre 多项式，
$a_p=(2p)!/[2^p(p!)^2]$ 为其最高次系数。
采用~\cite{vincent2011}~的归一化：
\[
\eta_p=\frac{c(2p+1)(a_pp!)^2}{2},
\]
\begin{equation}\label{eq:vcjh-functions}
\begin{aligned}
g_L(\xi)&=\frac{(-1)^p}{2}
\left[P_p(\xi)-\frac{\eta_pP_{p-1}(\xi)+P_{p+1}(\xi)}{1+\eta_p}\right],\\
g_R(\xi)&=\frac12
\left[P_p(\xi)+\frac{\eta_pP_{p-1}(\xi)+P_{p+1}(\xi)}{1+\eta_p}\right].
\end{aligned}
\end{equation}
$P_k(1)=1$、$P_k(-1)=(-1)^k$ 直接给出四个 FR 端点条件。
$c$ 为无量纲参数；改变参数时，界面通量仍保持一致，内部模式的提升随之变化。

\conceptbridge{为何引入含高阶导数的修正能量，而不是只检查端点条件？}
{端点条件只规定通量积分。能量估计还需控制内部模式与界面提升的配对。对最高阶模式调整范数，可以寻找使离散分部积分成立的提升算子；正定性保证该估计实际控制全部自由度。}
{守恒端点条件 $\longrightarrow$ 内部提升 $\longrightarrow$ 正定修正能量。}

\section{修正范数与稳定参数区间}
在精确积分的多项式基中，$D$ 为参考导数矩阵，$M$ 为质量矩阵，定义
\begin{equation}\label{eq:vcjh-norm}
K_c=\frac c2(D^p)^{\mathsf T}MD^p,\qquad H_c=M+K_c.
\end{equation}
于是
\[
U^{\mathsf T}H_cU
=\int_{-1}^1u_h^2\dd\xi
+\frac c2\int_{-1}^1(\partial_\xi^pu_h)^2\dd\xi.
\]
在 Legendre 模态基下，仅最高阶对角元变化：
\[
H_c=\diag\left(2,\frac23,\ldots,
\frac{2}{2p+1}+c(a_pp!)^2\right).
\]
因此
\begin{equation}\label{eq:vcjh-bound}
c>c_-:=-\frac{2}{(2p+1)(a_pp!)^2}
\quad\Longleftrightarrow\quad H_c>0
\quad\Longleftrightarrow\quad\eta_p>-1.
\end{equation}
因为 $D^{p+1}=0$，有 $K_cD=0$。修正函数~\eqref{eq:vcjh-functions}~的导数矩阵满足
\[
C=H_c^{-1}R^{\mathsf T}B,\qquad
H_cD+D^{\mathsf T}H_c=R^{\mathsf T}BR.
\]
这一矩阵解释与 SBP--CPR 构造相容~\cite{ranocha2016}。

\section{线性能量证明}
对常数 $a$，FR 方程乘 $U_j^{\mathsf T}H_c$。
离散分部积分将体积贡献化为两端迹的平方差，提升项补入共同通量。
周期求和得到
\[
\frac{\dd}{\dd t}\sum_j\frac{J_j}{2}U_j^{\mathsf T}H_cU_j
=\sum_f(\widehat f-a\{u_h\})[u_h]
=-\frac{|a|}{2}\sum_f[u_h]^2.
\]
常数模态满足 $K_c\bm1=0$，因此该提升同时保持物理单元平均。
负参数仍可在~\eqref{eq:vcjh-bound}~内给出正定范数；
接近下界时最高阶模式的权重变小，参数评价还需关注范数等价常数、谱半径和误差。

\section{特定参数关系与实际选择}
在本文的精确质量矩阵和线性平流设定下：
\begin{center}\small
\begin{tabularx}{\textwidth}{@{}l>{\raggedright\arraybackslash}X@{}}
\toprule
参数 & 具体关系\\
\midrule
$c=0$ & $H_c=M$，恢复第~\ref{ch:11}~章的精确积分 DG 提升。\\
$\eta_p=p/(p+1)$ & 对应内部通量点为 $P_p$ 零点、两端包含接口的特定一维 SD 通量构造。\\
$\eta_p=(p+1)/p$ & 得到 Huynh 的 $g_2$ 修正函数；与特定 Lobatto 求积 DGSEM 的关系需同时固定求积和节点。\\
\bottomrule
\end{tabularx}
\end{center}
这些关系以指定的修正函数、节点和求积为基础。
改变 SD 节点、非线性通量插值或空间维数时，可沿相同的体积与提升矩阵逐项比较。
文献中的参数若记作 $\kappa$ 且定义 $K=\kappa(D^p)^{\mathsf T}MD^p$，
则本章 $c=2\kappa$。
选择参数可依目标波数的色散、可用 CFL 和误差阈值比较；线性能量允许区间提供可分析的起点。

\aiexercise{以 $p=2$ 的 Legendre 基构造 $D,M,R,B$，对 $c=0$、SD 参数及一个负的稳定参数计算 $H_c,C$。检验正定性、$H_cC=R^{\mathsf T}B$ 和修正能量导数。}
{固定总自由度、上风通量及 RK4，分别测量上述参数的平滑波误差和稳定步长。报告半离散与全离散结果，并把非线性稳定性问题列为进一步阅读。}

\chapter{Mortar 方法与非匹配网格耦合}\label{ch:13}
\begin{chaptersummary}
局部高阶空间允许不同单元独立选取尺寸和次数，界面两侧的节点与迹空间因而可能不匹配。Mortar 在共同接口上组织投影、通量和回传，使两侧使用同一份积分数据。本章给出投影矩阵、常数保持及线性能量条件，为 hp 自适应和异构方法耦合提供接口。
\end{chaptersummary}
\learninggoals{建立共同接口分区和 Mortar 空间；计算迹投影与对偶回传；证明两侧积分通量抵消；说明能量稳定所需的求积与迹再现条件。}
\studyroute{略去曲面几何不一致、非线性熵投影和降阶 Mortar 的专门稳定设计；线性证明使用能再现两侧迹的共同空间。}{mortar method; nonconforming interface; mass-adjoint projection; hp nonconforming DG}{在同一线段上以左侧一次、右侧常数迹构造 Mortar 质量矩阵和投影，显式核对通量载荷与常数测试。}

\section{共同接口与迹空间}\index{Mortar 方法}
考虑两个二维单元相邻的一条平面边 $\Gamma$，法向固定从左到右。
若两侧次数不同或一条粗边对应多条细边，应先把物理接口划分为共同子面。
在每个子面上选取 Mortar 空间 $V_m=\Poly_{p_m}$，并用相同几何映射、法向和积分测度。
取 $p_m$ 至少为两侧迹次数，可在仿射子面上再现两侧迹。
该守恒耦合思路可见~\cite{kopriva1996}。

\conceptbridge{直接在两侧各自的节点上计算通量，为什么还需要共同接口？}
{节点集合不同会产生不同的通量采样和积分。共同接口将两侧迹转换到同一空间，
只计算一份数值通量，再以相反法向传回。常数保持保证总量一致，
相容的积分与对偶回传进一步连接两侧的能量配对。}
{非匹配迹 $\longrightarrow$ 共同接口空间 $\longrightarrow$ 共享通量与相容回传。}

\section{投影、通量与对偶回传}
设 Mortar 基为 $\mu_\alpha$，一侧迹基为 $\phi_i^s$，$s=L,R$，定义
\[
(M_m)_{\alpha\beta}=\int_\Gamma\mu_\alpha\mu_\beta\dd S,\qquad
(B_{ms})_{\alpha i}=\int_\Gamma\mu_\alpha\phi_i^s\dd S.
\]
$L^2$ 迹投影满足
\begin{equation}\label{eq:mortar-proj}
P_s=M_m^{-1}B_{ms},\qquad u_{s,m}=P_su_s.
\end{equation}
在共同积分节点以 $u_{L,m},u_{R,m}$ 评价单值法向通量；
用同一规则表示其载荷为 $M_mf_m^*$，则两侧的正向外通量载荷为
\begin{equation}\label{eq:mortar-return}
r_L=P_L^{\mathsf T}M_mf_m^*,\qquad
r_R=-P_R^{\mathsf T}M_mf_m^*.
\end{equation}
这些是弱残差的对偶载荷。若需要一侧的通量系数，可再作用该侧迹质量矩阵的逆；
该表达同时说明 $M_sP_{m\to s}=P_s^{\mathsf T}M_m$ 的质量伴随关系。

\section{守恒与线性能量条件}
设 $\bm1_s$ 和 $\bm1_m$ 表示常数函数。投影再现常数即
$P_s\bm1_s=\bm1_m$，于是
\[
\bm1_L^{\mathsf T}r_L+\bm1_R^{\mathsf T}r_R=0.
\]
一条粗边对应多个 Mortar 子面时，对这些子面的载荷求和；
粗细两侧仍使用同一组积分。

为得到线性能量估计，进一步假设共同空间精确再现两侧迹，
并且求积精确表示相应乘积。此时
$P_s^{\mathsf T}M_mP_s=M_s$，
若有多个子面则对左端矩阵求和。
对常法向速度 $a_n$ 取
\[
f_m^*=a_n\{u_m\}-\frac{|a_n|}{2}[u_m],
\]
两侧 DG 边界能量之和为
\[
-\frac{|a_n|}{2}[u_m]^{\mathsf T}M_m[u_m]\le0.
\]
若共同空间降阶、两侧几何或求积不一致，上述迹范数恒等式需要相应替代设计；
常数保持负责总量守恒，迹再现、求积与质量伴随关系进一步支撑能量估计。

\section{一个可手算的非匹配接口}
令 $\Gamma=[-1,1]$，左迹空间基为 $(1,s)$，右迹仅含常数；
Mortar 也取 $(1,s)$，精确积分得到
\[
M_m=\begin{pmatrix}2&0\\0&2/3\end{pmatrix},\quad
P_L=I_2,\quad P_R=\begin{pmatrix}1\\0\end{pmatrix}.
\]
若共同通量系数为 $(F_0,F_1)^{\mathsf T}$，则
\[
r_L=\begin{pmatrix}2F_0\\2F_1/3\end{pmatrix},\qquad r_R=-2F_0.
\]
常数测试只读取左侧第一分量，其总通量恰好抵消；
线性分量是左侧内部迹的残差。hp 加密改变接口空间时，需同步更新投影和物理测度。

\aiexercise{将例子改为左侧二次、右侧一次多项式，采用三个 Gauss 点精确积分，显式构造 $M_m,B_{mL},B_{mR}$。检查常数再现、质量伴随与上风界面能量。}
{把右侧边划分为两个子面并各用一次迹。写出子面映射的 Jacobian，证明粗侧总载荷等于两份细侧载荷的相反数。}

\chapter{现代高阶方法的组合与时间推进}\label{ch:14}
\begin{chaptersummary}
DG、FR 和 Mortar 已确定局部空间与接口算子，完整演化还需在时间上协调它们。本章从共享通量连接不同空间方法，再用 ADER 与 CERK 提供步内状态、LTS 协调局部步长、IMEX 处理刚性分拆。空间逼近、接口积分与时间推进由此组成共同的精度与守恒要求。
\end{chaptersummary}
\learninggoals{构造 DG--FR 的共享通量接口；说明 ADER、CERK、LTS 的不同作用；写出 IMEX 分拆；设计守恒的 h、p 与 hp 转移及时间收敛实验。}
\studyroute{略去高阶 RK 系数推导、时空预测器的一般非线性迭代及自适应误差估计证明；各具体算法需注明采用的积分和时间表。}{ADER DG; continuous explicit Runge--Kutta; conservative local time stepping; IMEX; hp adaptivity}{对相邻 DG 与 FR 单元使用同一界面求解器，先统一步长验证守恒，再引入二比一 LTS，比较界面时间积分与总量误差。}

\section{共享通量连接不同空间方法}
对同一守恒律，DG 与 FR 都可将单元平均更新写成外通量积分。
若接口上选择相同状态、法向和单值通量，并通过各自相容的提升算子回传，
可将一个 DG 单元与一个 FR 单元连接。非匹配接口使用第~\ref{ch:13}~章的 Mortar。
耦合的总体阶数受两侧逼近、接口投影和求积精度共同限制；
两侧原有能量估计还需在共同接口配对下闭合。

\conceptbridge{空间方法分别稳定后，为何组合时还要检查时间与接口？}
{稳定性估计包含界面通量与时间更新。组合方法改变邻居数据的时刻、投影和配对，
会改变这些贡献。共享通量积分控制总量平衡，稳定域和相容能量配对控制演化，
两类检查都需作用于最终组合算子。}
{局部空间算子 $\longrightarrow$ 接口与时间耦合 $\longrightarrow$ 全离散结构。}
\section{高阶时间推进：ADER、CERK 与 LTS}\label{sec:v1-time}
空间离散给出有限维系统
\[
M_h\dot U_h=R_h(U_h,t).
\]
对固定且可逆的质量矩阵，可写成 $\dot U_h=L_h(U_h,t)$。
时间格式近似这一系统的演化；其阶数、稳定域和界面时间数据必须与空间离散匹配。
以下分别讨论步内近似和局部时间调度，第二卷第 5.5 章给出 Maxwell 背景下的展开。

\paragraph{ADER：局部时空预测与守恒更新。}
ADER（Arbitrary high-order DERivatives）从高阶时间展开发展而来：
经典构造利用 PDE 将时间导数转为空间导数，现代构造也可在每个单元内求解局部时空弱问题。
由此得到 $q_K(\bm x,t)$，供体积分和界面数值通量在整个时间步上积分。
以无源有限体积更新为例，
\[
|K|(\bar U_K^{n+1}-\bar U_K^n)
=-\sum_{f\subset\partial K}\int_{t_n}^{t_{n+1}}\int_f
\widehat F(q_K,q_{K_f};\bm n_f)\dd S\dd t.
\]
单元局部预测之后，校正步骤通过共同通量耦合邻居。精度取决于空间重构、预测器与积分阶数；
显式双曲更新仍受相应 CFL 条件约束。
可参见 \href{https://doi.org/10.1016/j.cma.2014.07.019}{Dumbser 的 ADER--LTS 构造}。

在 ADER--DG 中，同一局部预测 $q_K$ 还提供基函数加权的体积残差。
将 DG 弱形式在一个时间步上积分，得到
\[
\begin{aligned}
\int_Kv_h(u_h^{n+1}-u_h^n)\dd x
={}&\int_{t_n}^{t_{n+1}}\int_K\nabla v_h\cdot F(q_K)\dd x\dd t\\
&-\int_{t_n}^{t_{n+1}}\int_{\partial K}v_h\widehat F(q_K,q_{K_f};n)\dd S\dd t.
\end{aligned}
\]
局部预测器给出时空数据，DG 测试与共享通量组成校正；这明确了 ADER 与空间 DG 的耦合位置。

\paragraph{CERK：连续显式 Runge--Kutta。}
CERK 在显式 RK 阶段导数 $k_i$ 的基础上构造步内连续近似
\[
U_h(t_n+\theta\Delta t)\approx
U_h^n+\Delta t\sum_{i=1}^{s}b_i(\theta)k_i,\qquad 0\le\theta\le1.
\]
其中 $b_i(0)=0$，$b_i(1)=b_i$ 与端点 RK 权重一致。
连续扩展的阶条件控制步内精度，端点 RK 的阶条件控制步末精度；两者分别满足相应的耦合与输出要求。
该近似可用于输出、耦合时间插值或局部预测器，其构造见
\href{https://doi.org/10.1137/0913084}{Owren--Zennaro 的连续显式 RK 方法}。
用于单元局部预测时，RK 阶段按局部问题组织，邻居影响由校正通量接入；
这种空间时间耦合可见~\cite{altmann2010}。

\paragraph{LTS：局部时间步进。}
若局部网格尺寸或波速差异较大，统一步长由最严格的单元约束决定。
LTS（local time stepping）令不同单元采用不同物理时间步长，并在共同时间层同步。
粗单元一次走过 $\Delta T$、细单元走过 $m$ 个子步时，共享面的时间积分通量须满足
\[
\Phi_f^{\mathrm{coarse}}
=\sum_{r=1}^{m}\Phi_{f,r}^{\mathrm{fine}},
\]
其中各项使用同一法向约定，在相邻单元更新中以相反符号出现。
ADER 或 CERK 可提供步内状态，LTS 还需配套调度、积分及通量累积。
真实时间 LTS 应保持时间精度；稳态加速中的局部伪时间步并不承担这一要求。
效率取决于小单元占比、界面预测、积分与同步成本。

\begin{askai}
\textbf{AI 展开问题：}\par
\textbf{略去内容与关键词：}ADER predictor--corrector、CERK dense-output order conditions、conservative LTS；这里略去高阶系数求解和调度算法。

对一维平流方程设置相邻粗细单元，令步长比为 $2:1$。分别用时间多项式和 CERK 连续扩展提供界面状态，列出两个细子步的通量积分，并验证粗细单元的总守恒量。

进一步展开：将粗侧状态在整步内冻结，测量时间收敛阶如何变化；再比较统一步长与 LTS 的更新次数。
\end{askai}

\section{RK 稳定域与 IMEX 分拆}
若空间物理模态满足 $\dot y=\lambda_hy$，RK 更新为
$y^{n+1}=\mathcal R(\Delta t\lambda_h)y^n$，
需将空间谱与时间稳定域共同检查。
SSP 方法将更新写成适当向前 Euler 步的凸组合，
其强稳定性继承需要 Euler 步在目标泛函下已有稳定性及相应步长约束。
空间高阶能量估计与 SSP 性质各有适用条件。

对传播和刚性松弛并存的问题，采用
\[
\dot U=F_h(U)+G_h(U)
\]
的 IMEX 分拆，$F_h$ 显式、$G_h$ 隐式。
自治加性 Runge--Kutta 的共同阶段写为
\[
\begin{aligned}
Y_i&=U^n+\Delta t\sum_{j<i}\widetilde a_{ij}F_h(Y_j)
+\Delta t\sum_{j\le i}a_{ij}G_h(Y_j),\\
U^{n+1}&=U^n+\Delta t\sum_i\widetilde b_iF_h(Y_i)
+\Delta t\sum_i b_iG_h(Y_i).
\end{aligned}
\]
两个系数表须满足加性阶条件~\cite{ascher1997}。
局部碰撞源的隐式处理可以减少快速松弛的步长限制；
显式传播仍需相应稳定域，刚性极限还需检验渐近保持性与实际时间阶数。
第二卷以导电 Maxwell 方程给出具体示例。

\section{h、p 与 hp 自适应}
$h$ 自适应改变单元尺寸，$p$ 自适应改变局部次数，$hp$ 自适应同时改变两者。
光滑区域可增加次数，局部间断和几何尺度常需缩小单元并配合稳定化。
转移算子应满足
\[
\int_Ku_{\mathrm{new}}\dd x
=\int_Ku_{\mathrm{old}}\dd x,
\]
细化时对全部子单元求和，粗化时先汇总子单元积分。
$L^2$ 投影含常数测试函数时保持该平均；正性及其他约束可能还需额外限制。
重分区后的接口空间、几何测度与积分规则必须一起更新。

\aiexercise{对周期平滑平流，固定空间网格，分别用 RK4 与 Heun 连续扩展测时间收敛，再固定足够小步长测空间收敛。用质量、能量和相位误差说明各指标反映什么。}
{对带快速 BGK 松弛的离散速度方程选一组有出处的 IMEX 系数，区分传播 CFL、碰撞松弛与未解析极限。将系数出处和实际收敛阶记录为实验的一部分。}

\part{介观动力学与 LBM}
\chapter{从 Boltzmann 方程到格子 Boltzmann 方法}\label{ch:15}
\begin{chaptersummary}
前面直接离散宏观方程，本章转而在动理学层次选择有限速度与碰撞模型，再通过长波、低 Mach 极限恢复宏观输运。D2Q9 与 D3Q27 将速度集合、各向同性矩和精确格点迁移联系起来，Chapman--Enskog 展开则说明非平衡应力与离散时间修正如何共同确定黏度。
\end{chaptersummary}
\learninggoals{定义离散粒子群与守恒矩；构造 D2Q9 与 D3Q27；完成迁移和碰撞更新；由多尺度展开解释低 Mach 流体极限和黏度公式。}
\studyroute{略去完整三阶 Chapman--Enskog 代数及一般热力学格子的构造；正文限于无外力、标准等温低 Mach 模型。}{lattice Boltzmann; D2Q9; D3Q27; Chapman--Enskog expansion; lattice viscosity}{验证 D2Q9 的二阶与四阶各向同性矩，再从迁移—碰撞方程的二阶 Taylor 展开推导黏度中的半时间步修正。}

\section{离散速度与介观变量}\index{LBM}
连续动理学含 $\bm x,\bm v$ 两类变量。
离散速度模型选 $\bm c_i$ 与粒子群密度 $f_i(\bm x,t)$，满足
\[
\partial_tf_i+\bm c_i\cdot\nabla f_i=\Omega_i,\qquad
\rho=\sum_i f_i,\qquad \rho\bm u=\sum_i\bm c_i f_i.
\]
本章的 $f_i$ 为质量密度型离散速度变量，已吸收速度求积权重，
区别于第一章的数密度型连续分布 $f$。
速度集合与平衡态的矩匹配决定可描述的宏观模型；
空间迁移与时间离散随后形成数值更新。
当 $\bm c_i\Delta t$ 恰好连接格点，特征迁移可直接实现为格点数据移位。

\conceptbridge{为何只保存有限个速度方向，还能得到宏观黏性流动？}
{宏观方程读取分布的低阶矩。若离散速度、权重和平衡态匹配所需的各向同性矩，
碰撞又保持质量与动量，长波展开便可计算非平衡应力。
恢复的方程与误差由这些矩条件、尺度排序和时间离散共同决定。}
{离散速度与碰撞矩 $\longrightarrow$ 非平衡应力 $\longrightarrow$ 宏观输运极限。}

\section{D2Q9 与 D3Q27}
令 $c_\ell=\Delta x/\Delta t$，D2Q9 的速度为
\[
\bm c_i/c_\ell\in\{(0,0),(\pm1,0),(0,\pm1),(\pm1,\pm1)\}.
\]
静止、轴向和对角方向的权重分别为
$4/9,1/9,1/36$，格子声速 $c_s^2=c_\ell^2/3$。
D3Q27 使用三个方向 $\{0,\pm c_\ell\}$ 的张量积，
一维权重 $w(0)=2/3$、$w(\pm c_\ell)=1/6$，三维权重为各方向权重的乘积。
两种格子都满足相应维数下的
\[
\sum_iw_i=1,\quad \sum_iw_i c_{i\alpha}=0,\quad
\sum_iw_i c_{i\alpha}c_{i\beta}=c_s^2\delta_{\alpha\beta},
\]
\[
\sum_iw_i c_{i\alpha}c_{i\beta}c_{i\gamma}c_{i\delta}
=c_s^4(\delta_{\alpha\beta}\delta_{\gamma\delta}
+\delta_{\alpha\gamma}\delta_{\beta\delta}
+\delta_{\alpha\delta}\delta_{\beta\gamma}).
\]
更高阶速度矩受到有限格子的限制。

\section{平衡态、碰撞与迁移}
等温二阶低 Mach 平衡态取
\begin{equation}\label{eq:lbm-eq}
f_i^{\mathrm{eq}}=w_i\rho\left[
1+\frac{\bm c_i\cdot\bm u}{c_s^2}
+\frac{(\bm c_i\cdot\bm u)^2}{2c_s^4}
-\frac{|\bm u|^2}{2c_s^2}\right].
\end{equation}
由上述矩恒等式，
\[
\sum_if_i^{\mathrm{eq}}=\rho,\quad
\sum_i\bm c_if_i^{\mathrm{eq}}=\rho\bm u,\quad
\sum_i\bm c_i\bm c_if_i^{\mathrm{eq}}=\rho c_s^2I+\rho\bm u\bm u.
\]
BGK--LBM 每步先作局部碰撞，再作迁移：
\begin{equation}\label{eq:lbe}
f_i^*(\bm x,t)=f_i-\frac{\Delta t}{\tau_{\mathrm{LB}}}(f_i-f_i^{\mathrm{eq}}),
\qquad
f_i(\bm x+\bm c_i\Delta t,t+\Delta t)=f_i^*(\bm x,t).
\end{equation}
碰撞保持零、一阶矩，周期格点迁移保持全局粒子群之和。
$\tau_{\mathrm{LB}}$ 有时间量纲；常用格子单位 $\Delta t=c_\ell=1$ 时，
松弛率为 $s=1/\tau_{\mathrm{LB}}$。边界反射或重构也需检查质量与动量交换。

\section{Chapman--Enskog 展开与黏度}\index{Chapman--Enskog 展开}
记 $D_i=\partial_t+\bm c_i\cdot\nabla$。
对式~\eqref{eq:lbe}~作光滑长波 Taylor 展开：
\[
D_if_i+\frac{\Delta t}{2}D_i^2f_i+\cdots
=-\frac{f_i-f_i^{\mathrm{eq}}}{\tau_{\mathrm{LB}}}.
\]
令空间与时间导数按小参数排序，并写
$f_i=f_i^{(0)}+\epsilon f_i^{(1)}+\epsilon^2f_i^{(2)}+\cdots$，
对非平衡部分施加零质量、零动量矩。
零阶给 $f_i^{(0)}=f_i^{\mathrm{eq}}$，一阶给
\[
f_i^{(1)}=-\tau_{\mathrm{LB}}D_i^{(1)}f_i^{\mathrm{eq}}.
\]
取零、一阶矩得到等温 Euler 主部；
二阶矩将 $f_i^{(1)}$ 转为黏性应力，Taylor 项进一步产生时间修正。
在低 Mach、近不可压缩的黏性尺度下，
\[
\partial_t\rho+\nabla\cdot(\rho\bm u)=0,
\]
\[
\rho_0(\partial_t\bm u+\bm u\cdot\nabla\bm u)
=-\nabla p+\rho_0\nu\Delta\bm u,\qquad
\nabla\cdot\bm u=0
\]
为相应极限方程；其中 $p$ 来自 $c_s^2\rho$ 的小幅压力扰动，
$\rho/\rho_0=1+O(\mathrm{Ma}^2)$，且
\begin{equation}\label{eq:lbm-viscosity}
\nu=c_s^2\left(\tau_{\mathrm{LB}}-\frac{\Delta t}{2}\right).
\end{equation}
$\tau_{\mathrm{LB}}>\Delta t/2$ 给出正黏度。
黏度正性给出松弛参数的基本范围，完整更新的稳定性还包括平衡态正性、边界与非流体模态。
连续 BGK 碰撞时间与离散 $\tau_{\mathrm{LB}}$ 的对应取决于时间积分；
式~\eqref{eq:lbm-viscosity}~适用于这里给定的更新。
标准等温二阶模型的宏观结论受低 Mach 和长波假设约束~\cite{qian1992}。

\aiexercise{列出 D2Q9 的九个速度与权重，直接计算平衡态的质量、动量和二阶矩。写出无边界周期格子的碰撞与迁移，逐步证明总质量守恒。}
{固定目标运动黏度，改变 $\Delta t$ 后重新计算 $\tau_{\mathrm{LB}}$；对周期剪切波 $u_x(y,0)=U_0\sin ky$ 测量衰减率 $\nu k^2$，区分格子单位与物理单位。}

\chapter{LBM 的现代模型与算法}\label{ch:16}
\begin{chaptersummary}
单松弛等温 LBM 已给出离散速度到宏观输运的路线。本章通过矩、累积量和离散熵比较不同碰撞设计，再讨论热、可压缩与边界模型需要增加的自由度和矩条件。模型选择由目标物理、Mach 数、温度变化与可验证的稳定范围共同决定。
\end{chaptersummary}
\learninggoals{区分原始矩、中心矩和累积量；写出 MRT 与离散熵约束；解释热与可压缩模型的额外矩需求；设计边界和 Galilean 误差检验。}
\studyroute{给出模型结构和基本公式，略去具体 MRT 矩阵表、累积量高阶回变换和特定可压缩格子的系数；这些内容按所选原始模型展开。}{MRT lattice Boltzmann; cumulant collision; entropic lattice Boltzmann; thermal lattice; Galilean invariance}{比较 BGK、MRT 与累积量碰撞保持哪些矩、分别调节哪些非守恒模式，再为同一剪切波规定统一黏度与误差指标。}

\section{MRT：在矩空间中调节松弛}\index{MRT--LBM}
对可逆矩变换 $m=Tf$，多松弛时间碰撞为
\[
m^*=m-S(m-m^{\mathrm{eq}}),\qquad
f^*=T^{-1}m^*.
\]
$S=\diag(s_a)$ 的守恒矩对应零松弛率，其余模式可有不同速率。
与剪切应力对应的速率 $s_\nu$ 决定
\[
\nu=c_s^2\Delta t\left(\frac1{s_\nu}-\frac12\right).
\]
其他非守恒模式的调节影响声学、体积黏性和格子尺度阻尼。
对局部对角松弛，$0<s_a<2$ 控制单独松弛模态；
完整更新还含迁移、平衡态线性化与边界，需检查组合谱~\cite{dhumieres2002}。

\conceptbridge{为何多松弛模型要指定所用矩，而不能只列出一组松弛率？}
{矩变换决定哪些模式对应守恒、应力和更高阶非平衡信息。
同一组速率作用在不同矩上会产生不同碰撞算子。
基于中心矩或累积量的碰撞还改变了变量与局部速度的关系。}
{碰撞坐标 $\longrightarrow$ 模态松弛 $\longrightarrow$ 输运与非流体误差。}

\section{中心矩与累积量}
原始矩为 $\sum_i f_i\bm c_i^\alpha$，中心矩以
$(\bm c_i-\bm u)^\alpha$ 替代。
若 $\rho>0$，用归一化粒子群 $f_i/\rho$ 定义生成函数
\[
K(\bm t)=\log\left(\frac1\rho\sum_i f_i
\exp[\bm t\cdot(\bm c_i-\bm u)]\right),\qquad
\kappa_\alpha=\left.\partial_{\bm t}^{\alpha}K\right|_{\bm t=0}.
\]
二阶累积量给出中心协方差，高阶累积量扣除了低阶矩乘积的贡献。
累积量碰撞在这些坐标中松弛，再重建 $f_i$，
从而调节不同阶非平衡信息的耦合~\cite{geier2015}。
有限速度集合仍限制可独立表达的矩，各模型的平衡态、速率与修正项需要一起指定。

\section{离散熵与正性}\index{熵格子 Boltzmann}
对正粒子群，可定义凸函数
\[
H(f)=\sum_i f_i\log\frac{f_i}{w_i}.
\]
熵平衡态在质量和动量约束下最小化 $H$。
一种熵碰撞沿 $\Delta=f^{\mathrm{eq}}-f$ 寻找保持正性的非零根
$H(f+\alpha\Delta)=H(f)$，再用 $0<\beta\le1$ 取
$f^*=f+\alpha\beta\Delta$。
凸性给出该线段上的 $H(f^*)\le H(f)$。
上述熵估计以熵平衡态、正性域及非零根存在为前提。采用多项式平衡态或固定 BGK 速率时，应对相应碰撞映射重新建立这些条件。
原始离散熵构造见~\cite{karlin1999}。

\section{热模型、可压缩模型与边界}
等温格子的 $p=c_s^2\rho$ 固定了相应状态关系。
热传输需增加能量相关的离散分布或更丰富速度集合，
并匹配热流矩、比热和目标 Prandtl 数。
可压缩流还需正确的能量通量、可变温度响应及足够阶的速度矩；
具体模型按速度集合、平衡展开与碰撞形式验证其宏观极限。

静止壁面的半程反弹将出射粒子群反向为入射粒子群；
壁面位置、滑移误差与几何处理相关。移动壁面需补动量交换，
入口与出口则需按给定宏观条件重建未知入射群。
热边界还需给定能量交换。曲面插值、漫反射与稀薄气体边界有各自适用条件。

\section{稳定性与 Galilean 误差的比较}
对剪切波叠加均匀背景速度 $U_b$，连续流体的衰减率应与 $U_b$ 无关。
比较不同背景速度下的黏度与相位，可测量格子高阶矩造成的 Galilean 误差。
同时记录 Mach 数、黏度、格子分辨率、边界位置及粒子群正性。
稳定性分析可结合线性谱、离散熵与非线性实验，不同指标针对不同误差机制。

\aiexercise{从 D2Q9 中选择一个明确的可逆矩变换，标出守恒与剪切模式。令非守恒速率相同，核对其在相应平衡矩下何时恢复 BGK 更新。}
{对周期剪切波添加两个不同背景速度，比较 BGK、MRT 和一个有出处的累积量模型的衰减误差；再说明若要计算温度变化还缺哪些状态量与矩约束。}

\part{复杂物理系统中的约化与闭合}
\chapter{稀薄气体动力学与非平衡闭合}\label{ch:17}
\begin{chaptersummary}
第一、二章通过速度矩把动理学联系到连续介质；当非平衡增强时，需要重新选择保留的状态与闭合。本章按 Knudsen 数比较 Chapman--Enskog、Grad 和简化碰撞模型，再以 DSMC 展示对 Boltzmann 动力学的随机粒子近似。由此可以分别追踪模型、闭合与采样误差。
\end{chaptersummary}
\learninggoals{区分 Chapman--Enskog 与有限矩闭合；解释 Grad 13 矩和 ES--BGK 的自由度；用局部尺度评估连续介质适用性；辨认 DSMC 采样误差。}
\studyroute{略去高阶矩系统的双曲性证明、一般 H 定理和 DSMC 碰撞采样的完整构造；输运系数按具体碰撞核计算。}{Grad 13 moment; Chapman--Enskog; ES BGK; Knudsen layer; DSMC}{用平面 Couette 流比较 Newton 应力、Grad 应力和动理学分布，指出 Knudsen 层中的边界数据和体域闭合分别缺少什么。}

\section{Knudsen 数与非平衡信息}\index{稀薄气体}
$\mathrm{Kn}=\lambda/L$ 比较平均自由程与宏观尺度。
常用的连续、滑移、过渡和自由分子区间是建模指南，
局部梯度、壁面层和目标精度还会改变具体模型的适用性。
例如 $L_T=T/|\nabla T|$ 可形成温度梯度 Knudsen 数 $\lambda/L_T$。
即使整体 $\lambda/L$ 很小，激波和壁面附近也可能需要非平衡描述。

\section{渐近流形与矩系统的两条闭合路径}
Chapman--Enskog 展开在小 Knudsen 数下写作
\[
f=M[n,\bm u,T]+\mathrm{Kn}\,f^{(1)}
+\mathrm{Kn}^2f^{(2)}+\cdots.
\]
非平衡修正的守恒矩取零，使 $n,\bm u,T$ 由 Maxwellian 参数定义。
第一修正经速度积分给出 Newton 应力与 Fourier 热流；
输运系数依赖线性化碰撞算子。

Grad 方法则直接保留更高矩作为独立演化量。
三维单原子气体的 13 矩由
$n$、$\bm u$、$T$、五个无迹应力分量与三个热流分量组成。
一种相应的分布表示结构为
\[
f_{13}=M\left[1+
A_{ij}\left(c_ic_j-\tfrac13|\bm c|^2\delta_{ij}\right)
+B_i c_i(|\bm c|^2-5\theta)\right],
\qquad \theta=k_BT/m.
\]
$A_{ij}$ 对称无迹，$A,B$ 由指定应力与热流确定。
将该表示代入更高矩，形成有限矩闭合~\cite{grad1949}。
多项式修正可能在速度尾部变负，系统的可实现性与双曲性需要专门检查；
正则化矩模型按其适用范围另行选择。

\conceptbridge{为何同样求应力和热流，渐近闭合与 Grad 闭合会给出不同模型？}
{渐近闭合以小 Knudsen 数约束非平衡分布，使应力与热流成为宏观梯度的函数。
Grad 模型将这些矩保留为独立变量，再补足其方程中的更高矩。
两者选择的状态空间和时间尺度不同，误差与适用区域也相应不同。}
{小参数渐近关系 / 独立高阶矩 $\longrightarrow$ 两类闭合模型。}

\section{BGK 与 ES--BGK}
连续 BGK 以 Maxwellian 为松弛目标，保持质量、动量和能量；
其单一松弛时间把不同输运过程联系起来，标准单原子模型给出 $\mathrm{Pr}=1$。
ES--BGK 使用各向异性高斯：
\[
G_\chi=\frac{n}{(2\pi)^{3/2}\sqrt{\det\mathcal T_\chi}}
\exp\left(-\frac12\bm c^{\mathsf T}\mathcal T_\chi^{-1}\bm c\right),
\]
\[
\mathcal T_\chi=(1-\chi)\theta I+\chi\Theta,\qquad
\Theta=\bm P/\rho,\qquad
C[f]=\frac{G_\chi-f}{\tau_c}.
\]
$\tr\Theta=3\theta$ 保持能量，协方差必须正定。
该参数化允许调节 $\mathrm{Pr}=1/(1-\chi)$；
单原子气体常用 $\chi=-1/2$ 对应 $\mathrm{Pr}=2/3$，
并需按该碰撞频率定义匹配黏度~\cite{holway1966}。
离散速度、求积或时间推进改变后，还应检查碰撞守恒矩及熵性质。

\section{DSMC 与误差分层}
DSMC 用带权模拟粒子表示分布，在小单元中交替自由迁移与随机碰撞。
碰撞对和散射角按 Boltzmann 碰撞核抽样；
边界入射粒子由壁面或入口分布产生~\cite{bird1994}。
有限粒子数引入统计波动，单元大小和时间步长带来数值偏差。
解析平均自由程、碰撞时间和局部梯度是相应精化目标。

比较连续介质与动理学模型时，可先固定碰撞物理，
再分别减小网格、步长和抽样误差。
这样才能将模型闭合误差与数值误差区分开。

\aiexercise{取标准正态密度乘以 $1+\alpha H_3(v)$，其中 $H_3(v)=v^3-3v$。说明任意非零 $\alpha$ 为何可能使速度尾部取负。再计算其零、一、二阶矩，解释低阶矩匹配与正性之间的关系。}
{设计一组 Couette 流比较：分别改变 Knudsen 数、速度空间分辨率和 DSMC 样本数，测量滑移速度与应力。区分闭合误差、边界模型误差和抽样误差。}

\chapter{MHD 与等离子体流体模型}\label{ch:18}
\begin{chaptersummary}
Maxwell 与物种动理学描述包含电荷分离、电子及离子时间尺度。MHD 利用低频、准中性与尺度排序形成单流体模型，再通过压力、热流和 Ohm 关系闭合。本章给出理想 MHD 的守恒形式，并以电阻和压力各向异性说明物理闭合如何改变能量与数值约束。
\end{chaptersummary}
\learninggoals{列出理想 MHD 的未知量与假设；核对总能量和磁散度；区分电阻率与磁扩散率；解释压力张量闭合及其适用尺度。}
\studyroute{略去从双流体到广义 Ohm 定律的完整渐近推导、输运系数表及各向异性不稳定性的色散计算。}{ideal MHD; resistive MHD; generalized Ohm law; gyrotropic pressure; CGL closure}{从理想 MHD 的动量、感应和内能方程推导总能量通量，逐项核对磁压与磁张力，并说明使用磁散度约束的位置。}

\section{低频单流体状态}\index{MHD}
本章使用 SI 单位，取 $\rho,\bm u,p,\bm B$ 为单流体变量。
准中性使主体电荷密度相对较小；低频排序下忽略位移电流，得到
\[
\bm J=\frac1{\mu_0}\nabla\times\bm B,\qquad
\partial_t\bm B=-\nabla\times\bm E.
\]
理想 Ohm 关系为 $\bm E+\bm u\times\bm B=0$。
其适用尺度还需使 Hall、电子压力梯度和电子惯性等项相对较小；
例如离子惯性长度与目标长度之比必须与所选排序相容。
压强取各向同性，并给定理想气体闭合 $p=(\gamma-1)\rho e$。

\section{理想 MHD 的守恒系统}
令
\[
\mathcal E=\rho e+\tfrac12\rho|\bm u|^2+\frac{|\bm B|^2}{2\mu_0},
\qquad p_t=p+\frac{|\bm B|^2}{2\mu_0}.
\]
无外力、黏性和热传导时，
\begin{equation}\label{eq:ideal-mhd}
\begin{aligned}
\partial_t\rho+\nabla\cdot(\rho\bm u)&=0,\\
\partial_t(\rho\bm u)+\nabla\cdot
\left(\rho\bm u\bm u+p_tI-\frac{\bm B\bm B}{\mu_0}\right)&=0,\\
\partial_t\bm B-\nabla\times(\bm u\times\bm B)&=0,\\
\partial_t\mathcal E+\nabla\cdot
\left[(\mathcal E+p_t)\bm u-\frac{(\bm u\cdot\bm B)\bm B}{\mu_0}\right]&=0,
\end{aligned}
\end{equation}
并要求 $\nabla\cdot\bm B=0$。
磁压与磁张力形式使用该约束；连续传播依赖 $\nabla\cdot\nabla\times=0$。
离散设计需同时验证单元守恒、磁散度、正密度与正内能。

\conceptbridge{为何写出守恒型 MHD 后，还要保留独立的磁散度与闭合条件？}
{守恒方程描述各状态量的通量交换；磁散度限定允许的磁场。
压力和 Ohm 关系则来自尺度与输运假设。三者分别约束状态、演化与未解物理，
并在离散后对应不同的验证指标。}
{低频约化 $\longrightarrow$ 单流体守恒系统 + 磁约束 + 物理闭合。}

\section{电阻、扩散与能量转换}
取电阻率 $\eta_e\ge0$，
\[
\bm E+\bm u\times\bm B=\eta_e\bm J.
\]
在 $\eta_e,\mu_0$ 为常数且 $\nabla\cdot\bm B=0$ 时，
\[
\partial_t\bm B=\nabla\times(\bm u\times\bm B)
+\eta_m\Delta\bm B,\qquad \eta_m=\eta_e/\mu_0.
\]
$\eta_e$ 的单位为 $\Omega\cdot\mathrm m$，$\eta_m$ 的单位为 $\mathrm m^2/\mathrm s$。
磁能的 Ohmic 损耗为 $\eta_e|\bm J|^2$，内能获得同样的焦耳热；
总能量还需使用与电阻电场相容的 Poynting 通量。
显式磁扩散通常给出 $O(h^2/\eta_m)$ 步长限制，可按传播与扩散尺度选择 IMEX。

\section{压力各向异性与输运闭合}
强磁场中可写
\[
\bm P=p_\perp I+(p_\parallel-p_\perp)\bm b\bm b+\bm\Pi_{\mathrm{ng}},
\qquad \bm b=\bm B/|\bm B|.
\]
$\bm\Pi_{\mathrm{ng}}$ 为非回旋对称部分；热流也可区分平行与垂直分量。
在碰撞可忽略、热流与有限回旋半径作用可忽略的双绝热近似中，
\[
\frac{D}{Dt}\left(\frac{p_\perp}{\rho|\bm B|}\right)=0,\qquad
\frac{D}{Dt}\left(\frac{p_\parallel|\bm B|^2}{\rho^3}\right)=0.
\]
CGL 关系将平行与垂直压力分别联系到密度和磁场，构成上述双绝热排序下的压力闭合。
碰撞输运或非局域热流需使用对应的模型，且可能改变适定性和不稳定性阈值~\cite{cgl1956}。

\aiexercise{在周期区域内，对常电阻率 MHD 分别计算磁能损耗与内能焦耳热，证明体积转换项相抵，并写出边界能量通量。}
{对平衡磁场叠加小幅 Alfv\'en 波，测波速、相位和离散磁散度。再引入平行与垂直压力，列出原有各向同性闭合中需要改变的方程。}

\chapter{回旋动理学与多尺度约化}\label{ch:19}
\begin{chaptersummary}
第 4 章的强磁场排序将快回旋与慢输运分离，本章据此构造回旋动理学。坐标变换与有序平均消去快角度，保留平行速度和磁矩分布；回旋平均与物理矩推回则重新连接约化粒子和自洽场。均匀磁场静电模型给出这条五维约化路线的具体形式。
\end{chaptersummary}
\learninggoals{列出强磁场排序与五维坐标；区分导引中心与回旋中心；计算 Fourier 回旋平均；说明动理学约化、矩闭合、场耦合和数值离散各自的步骤。}
\studyroute{正文给出受控的均匀磁场静电例子，略去一般电磁 Lie 变换、二阶 Hamiltonian、磁化项和磁几何碰撞变换；具体模型须按同一变分阶数配套。}{guiding center; gyrocenter; gyroaverage; finite Larmor radius; gyrokinetic polarization; electromagnetic gyrokinetics}{在均匀磁场中对一个垂直 Fourier 电势作回旋平均，推导 Bessel 因子，再由指定的非绝热分布定义写出线性准中性关系。}

\section{尺度排序与五维状态}\index{回旋动理学}
以平衡尺度 $L_B$、离子热回旋半径 $\rho_{Li}$ 为参照，一组常用微扰排序为
\[
\epsilon=\rho_{Li}/L_B\ll1,\qquad
\omega/\Omega_i=O(\epsilon),\qquad
k_\parallel/k_\perp=O(\epsilon),\qquad
k_\perp\rho_{Li}=O(1).
\]
弱扰动还可取 $|q_i\delta\phi|/T_i=O(\epsilon)$；
本章后文的 $T_s$ 表示能量单位温度，即前文的 $k_BT_s$。
具体电磁、碰撞和电子尺度模型还需补充相应排序。

原相空间 $(\bm x,\bm v)$ 可改写为
$(\bm R,v_\parallel,\mu,\vartheta)$，
其中 $\mu=m v_\perp^2/(2B_0)$，$\vartheta$ 是快回旋角。
先由导引中心变换处理平衡场的快回旋，
再由回旋中心变换处理扰动引起的角度依赖。
约化分布 $F_s(\bm R,v_\parallel,\mu,t)$ 不再独立依赖 $\vartheta$，
保留三个位置和两个速度变量。
在均匀 $B_0$ 下，对回旋角无关的函数 $A$，速度矩使用
\[
\int A\dd^3v=\frac{2\pi B_0}{m_s}
\int_{-\infty}^{\infty}\int_0^\infty A\dd\mu\dd v_\parallel.
\]
角度积分已解析完成，平行速度与磁矩仍保留为独立变量。

\conceptbridge{消去一个快角度后，场方程为何还需要极化和回旋平均？}
{位置坐标从粒子位置变为回旋中心，密度与电流仍须在物理位置测量。
把约化分布推回物理空间时，会出现有限回旋半径、极化与磁化贡献。
这些贡献与约化动力学来自同一阶坐标变换，决定自洽场及能量配对。}
{相空间坐标变换 $\longrightarrow$ 五维演化 $\longrightarrow$ 推回物理矩与自洽场。}

\section{回旋平均与有限回旋半径}
令粒子位置 $\bm x=\bm R+\bm\rho(\vartheta)$。
在固定回旋中心与磁矩下定义
\[
\langle\phi\rangle_{\bm R}
=\frac1{2\pi}\int_0^{2\pi}
\phi(\bm R+\bm\rho(\vartheta))\dd\vartheta.
\]
对均匀磁场中的平面波，
\[
\left\langle\ee^{\ii\bm k\cdot\bm x}\right\rangle_{\bm R}
=\ee^{\ii\bm k\cdot\bm R}
J_0(k_\perp v_\perp/\Omega_s).
\]
当 $k_\perp\rho_{Ls}$ 有限时，环平均改变粒子对场的响应。
小参数消去快角度，Bessel 因子则保留其对慢尺度耦合的影响。

\section{动力学方程的结构与一个具体模型}
一般回旋中心方程可按特征演化表示为
\[
\partial_tF_s+\dot{\bm R}_s\cdot\nabla_{\bm R}F_s
+\dot v_{\parallel s}\partial_{v_\parallel}F_s=C_s^{\mathrm{GK}}[F].
\]
特征、相空间测度和碰撞项需由同一约化坐标构造；
非均匀磁场下的守恒形式使用对应的相空间 Jacobian。

作为可计算例子，取均匀 $B_0\bm b$、无平衡梯度、无碰撞、静电低频扰动。
令 $F_{0s}$ 为温度 $T_s$ 的均匀 Maxwellian，
$h_s$ 为通常的非绝热分布变量，其规范由下节物理密度公式固定。
一组对应的非线性静电方程为
\begin{equation}\label{eq:slab-gk}
\partial_th_s+v_\parallel\partial_\parallel h_s
+\bm v_{E,s}\cdot\nabla_\perp h_s
=\frac{q_sF_{0s}}{T_s}\partial_t\langle\phi\rangle_{\bm R},
\qquad
\bm v_{E,s}=\frac{\bm b\times\nabla\langle\phi\rangle_{\bm R}}{B_0}.
\end{equation}
漂移速度等于环平均电场的 $\bm E\times\bm B/B_0^2$ 速度。
均匀平衡下不含梯度驱动；加入密度、温度梯度后需配套加入驱动项。
一般电磁构造与能量守恒见~\cite{brizard2007,abel2013}。

\section{准中性、极化与 Amp\`ere 耦合}
对上述均匀磁场模型的 Fourier 模态，物理密度扰动为
\[
\delta n_{s,\bm k}
=-\frac{q_sn_{0s}}{T_s}\phi_{\bm k}
+\int J_{0s}h_{s,\bm k}\dd^3v,\qquad
J_{0s}=J_0(k_\perp v_\perp/\Omega_s).
\]
准中性条件 $\sum_sq_s\delta n_s=0$ 因而给出
\begin{equation}\label{eq:gk-qn}
\sum_s\frac{q_s^2n_{0s}}{T_s}\phi_{\bm k}
=\sum_sq_s\int J_{0s}h_{s,\bm k}\dd^3v.
\end{equation}
若改用 $g_{s,\bm k}=h_{s,\bm k}-(q_sF_{0s}/T_s)J_{0s}\phi_{\bm k}$，
同一关系变成
\[
\sum_s\frac{q_s^2n_{0s}}{T_s}(1-\Gamma_{0s})\phi_{\bm k}
=\sum_sq_s\int J_{0s}g_{s,\bm k}\dd^3v,
\]
\[
\Gamma_{0s}=I_0(b_s)\ee^{-b_s},\qquad
b_s=k_\perp^2T_s/(m_s\Omega_s^2).
\]
两种写法使用不同分布变量，极化项的位置因此改变。
零波数的电势规范与粒子数兼容条件需要另行固定。

电磁模型还需由平行电流确定 $A_\parallel$，
并在相应排序中处理 $\delta B_\parallel$、极化和磁化。
Amp\`ere 方程中的电流必须使用与所选回旋中心坐标、分布变量和场方程一致的推回。
这些场关系连接保留的速度分布与电磁势；高阶矩截断属于另一种闭合选择。

\section{碰撞、湍流输运与离散层次}
回旋动理学碰撞模型需要保持适当的粒子、动量与能量矩，
并与约化测度和熵耗散相容。湍流粒子与热通量由分布扰动和径向漂移的相关量产生，
例如粒子通量具有 $\langle\int\delta f_s\,v_{E,r}\dd^3v\rangle$ 的结构。
具体通量需按推回定义和平均区域计算。

五维状态仍需在位置、平行速度、磁矩与时间上离散。
域范围、速度尾截断和碰撞分辨率分别控制不同误差；
选用谱、DG 或粒子表示属于随后发生的数值选择。

\aiexercise{对 Maxwellian 平衡和均匀磁场，计算 $\langle\exp(\ii\bm k_\perp\cdot\bm\rho)\rangle$，再验证 $\int F_{0s}J_{0s}^2\dd^3v=n_{0s}\Gamma_{0s}$。由此证明 $h_s$ 与 $g_s$ 两种准中性写法一致。}
{从式~\eqref{eq:slab-gk}~出发，固定周期边界和速度截断，列出验证准中性残差、自由能与速度空间收敛所需的数据。进一步阅读时选定一套电磁模型，逐项对应其分布变量、Jacobian、碰撞项与场方程。}

\part{计算物理中的网格与验证}
\chapter{网格是几何与离散算子的接口}\label{ch:20}
\begin{chaptersummary}
模型和离散空间确定后，网格将物理几何、边界和局部尺度引入算子。本章从参考单元映射出发，解释形状与尺寸如何改变梯度、积分及界面配对，再联系曲边几何和守恒检验。这些几何相容性要求也构成下一章复杂网格技术的共同基础。
\end{chaptersummary}
\learninggoals{计算参考到物理单元的梯度与积分变换；区分几何误差与逼近误差；说明网格形状、邻接及边界表示对空间算子的影响。}
\studyroute{略去高维网格生成、几何近似误差证明和矩阵实现；本章保留算子结构，代数求解与并行数据布局见第 0 卷。}{element mapping; mesh quality; curved element; metric identity; discrete geometric conservation}{比较等边三角形与细长三角形的 Jacobian 和线性有限元刚度矩阵，解释梯度放大来自哪一个几何因子。}
\section{几何、邻接与离散耦合}\index{网格几何}
网格决定几何表示误差、局部离散误差、稳定性限制和矩阵非零结构。几何处理与网格质量控制是计算工作流的重要环节。

\begin{authornote}
\textbf{潜在疑问：}
“为什么网格不仅影响几何外形，还会影响矩阵结构、稳定性限制和求解器行为？”

\textbf{作者补充：}
网格不仅给出采样位置，还规定单元之间如何交换信息。邻接关系决定离散耦合，单元形状和尺寸进入导数、积分与时间步长估计，边界表示影响边界条件施加。同一连续方程在不同网格上会产生不同的离散耦合和误差。

\textbf{核心概念联系：}
\[
\text{mesh geometry}
\longrightarrow
\text{discrete coupling}
\longrightarrow
\text{matrix and solver behavior}
\]
\end{authornote}

贴体网格使网格边界与物体边界一致。其边界条件施加自然，适合解析壁面边界层，但复杂几何和大运动会导致网格生成与质量控制困难。

\section{参考单元映射与几何尺度}
设 $x=F_K(\xi)$，$A_K=\partial F_K/\partial\xi$ 可逆，则
\[
\nabla_xu=A_K^{-\mathsf T}\nabla_\xi\widehat u,\qquad
\dd x=|\det A_K|\dd\xi.
\]
Poisson 刚度条目成为
\[
K_{ij}^{(K)}
=\int_{\widehat K}(\nabla_\xi\phi_i)^{\mathsf T}
A_K^{-1}A_K^{-\mathsf T}\nabla_\xi\phi_j\,|\det A_K|\dd\xi.
\]
细长或退化单元改变映射奇异值，放大某些方向的梯度及误差常数。
结构化网格具有规则连接；非结构网格允许更灵活的局部形状与邻接。
两者都需规定边界取向、面测度和离散积分。

\section{曲边单元与几何误差}
高阶场近似常需相应阶数的曲边几何。
参考到物理映射的次数、体积求积和面求积共同影响总误差。
在守恒律中，常数解保持要求离散几何因子满足对应的度量恒等式；
共享面几何和相反法向是通量抵消的基础。
验证时可先用常数场、低阶多项式和解析曲面检查几何积分，再评估物理解的收敛。

\aiexercise{对映射 $x=A\xi+b$，分别取旋转、各向同性缩放和强拉伸的 $A$。推导其梯度、质量与刚度缩放，并比较三种网格上的同一线性场误差。}
{将一个圆边界分别用折线和二次曲线表示，固定场近似阶数，设计区分几何误差与空间逼近误差的收敛实验。}

\chapter{复杂几何与网格技术}\label{ch:21}
\begin{chaptersummary}
固定贴体网格直接表示物体边界，复杂运动与局部尺度变化则要求更灵活的几何组织。本章依次讨论浸没边界、树网格、ALE 和重叠网格，说明它们如何通过附加作用、共同子面、移动控制体与插值连接几何和演化。守恒转移与几何守恒律贯穿这些不同布局。
\end{chaptersummary}
\learninggoals{解释四类网格技术的状态与接口；由移动控制体推导相对通量；说明粗细转移、共享子面与重叠插值中的守恒要求。}
\studyroute{正文围绕几何表示、界面传递与守恒组织网格方程；网格生成器、孔切搜索及边界强迫的完整实现留待进一步展开。}{immersed boundary; quadtree octree; AMR refluxing; ALE; geometric conservation law; overset}{对一个粗边与两条细边显式列出共同通量，再对一个移动控制体检验常数解，比较空间与几何守恒的要求。}
\section{浸没边界与背景网格}
浸没边界法把复杂物体放入简单背景网格，不要求网格贴合物体。在动量方程中加入边界作用：
\[
\rho\left(
\frac{\partial\bm u}{\partial t}
+\bm u\cdot\nabla\bm u
\right)
=
-\nabla p+\mu\nabla^2\bm u+\bm f_{\mathrm{IB}}.
\]

\begin{authornote}
\textbf{潜在疑问：}
“为什么不让网格贴合边界时，可以通过方程里的 \(\bm f_{\mathrm{IB}}\) 表达边界作用？”

\textbf{作者补充：}
贴体网格把物体表面直接表示为网格边界，边界条件可以在该边界上施加。浸没边界法允许几何表面穿过背景网格，于是需要借助附加力或修改附近离散方程，使数值解满足物体的影响。几何处理更灵活，但边界附近的精度和守恒性需要专门检验。

\textbf{核心概念联系：}
\[
\text{body-fitted boundary condition}
\longrightarrow
\text{immersed forcing}
\longrightarrow
\text{background-grid computation}
\]
\end{authornote}

连续强迫法通过光滑核在拉格朗日边界与欧拉网格之间传递力；离散强迫法直接修改边界附近离散方程。

规则背景网格便于组织计算和处理复杂运动，壁面附近精度、通量守恒及薄结构则需要相应的边界离散与验证。

\section{四叉树、八叉树与 AMR}
二维四叉树将一个单元递归分成四个子单元，三维八叉树分成八个子单元。其核心是只在激波、剪切层、涡、界面或高误差区域投入高分辨率。

\begin{authornote}
\textbf{潜在疑问：}
“为什么局部加密会立刻带来 2:1 平衡、非匹配界面和守恒转移这些工程约束？”

\textbf{作者补充：}
局部加密改变的不只是分辨率：粗细单元之间会出现新的邻接关系，界面通量必须在不同尺度的面之间保持相容。随时间变化的加密区域还要求重新构造空间投影与界面通量，使新旧自由度之间的转移满足守恒。

\textbf{核心概念联系：}
\[
\text{local refinement}
\longrightarrow
\text{nonmatching interfaces}
\longrightarrow
\text{conservative transfer}
\]
\end{authornote}

常见工程约束包括：
\begin{itemize}
  \item 相邻单元保持 2:1 层级平衡；
  \item 处理粗细非匹配界面；
  \item 加密和粗化时守恒地转移变量；
  \item 自适应后重建接口投影与通量匹配。
\end{itemize}

\section{动网格、ALE 与几何守恒律}
在任意拉格朗日--欧拉描述中，网格具有速度 \(\bm w\)。守恒方程写成
\[
\frac{\dd}{\dd t}\int_{\Omega(t)}U\dd V
+
\int_{\partial\Omega(t)}
[\bm F(U)-U\bm w]\cdot\bm n\dd S
=0.
\]

\begin{authornote}
\textbf{潜在疑问：}
“为什么网格一运动，通量就从 \(\bm F(U)\) 变成 \(\bm F(U)-U\bm w\)？”

\textbf{作者补充：}
控制体运动时，守恒量穿过边界的变化不仅来自物质相对固定空间的运动，也来自边界自身扫过的体积。积分守恒律因此涉及物理通量与网格速度共同构成的相对通量。几何守恒律要求离散体积变化与网格扫过量一致，否则连均匀解也可能产生人为变化。

\textbf{核心概念联系：}
\[
\text{fixed control volume}
\longrightarrow
\text{moving control volume}
\longrightarrow
\text{ALE flux}
\]
\end{authornote}

动网格必须满足几何守恒律：若物理解为常数，仅仅移动网格不应制造虚假变化。

网格运动可由弹簧类比、Laplace 平滑、线弹性、径向基函数或局部重构实现。运动过大时，需要重新划分网格。

\section{重叠网格与非匹配耦合}
重叠网格将多个独立网格块叠加，例如背景网格、贴体部件网格和随动局部网格。网格之间通过供体单元向接收单元插值。

\begin{authornote}
\textbf{潜在疑问：}
“为什么多个网格叠在一起仍能形成一个计算域？插值耦合会牺牲什么结构？”

\textbf{作者补充：}
重叠网格用多个易于生成或运动的子网格共同覆盖计算区域。为避免重复求解，需标记被遮盖的单元，并在网格交界处从 donor 单元向 receiver 提供插值数据。插值建立了子网格之间的耦合，但普通插值不一定使两侧通量精确平衡，因此全局守恒性需要额外设计或检查。

\textbf{核心概念联系：}
\[
\text{single conforming mesh}
\longrightarrow
\text{overlapping mesh blocks}
\longrightarrow
\text{interpolation coupling}
\]
\end{authornote}

其优势是部件网格可独立生成，适合多个刚体的大范围相对运动。困难在于孔切、供体搜索、守恒插值和动态连通更新。


\section{空间、时间与几何的一致性}\index{几何守恒律}
粗细接口可使用 Mortar 组织共同子面，重叠网格则还需明确物理域的有效覆盖，
避免重复计量守恒量。普通状态插值、通量匹配与守恒重映射分别承担不同任务。

对移动单元，连续几何恒等式为
\[
\frac{\dd |K(t)|}{\dd t}
=\int_{\partial K(t)}\bm w\cdot\bm n\dd S.
\]
离散体积变化与 ALE 通量中的网格扫过体积需采用一致时间积分。
将 $U=U_0$ 代入移动控制体更新，即可检验均匀场是否保持。

\aiexercise{令一个一维单元的左右端点分别以给定速度运动，从 ALE 积分式推导常数解更新，写出体积时间积分与网格通量必须满足的等式。}
{画出树网格粗细交界与两块重叠网格，分别标出状态转移、有效控制体和共享通量。说明每种布局如何测量整体质量变化。}

\chapter{方法选择与完整计算工作流}\label{ch:22}
\begin{chaptersummary}
从状态表示、尺度约化与闭合，到空间离散和网格，全书已形成完整的建模链。本章以目标物理量统筹这些选择，将模型、约化、空间、时间与代数误差分层检验，再把守恒、稳定和收敛组织为可执行流程。最终方法组合由目标量误差、结构保持和验证证据共同确定。
\end{chaptersummary}
\learninggoals{为给定物理目标选择模型和离散；区分五类误差；设计空间、时间和容差独立变化实验；区分数值验证与物理确认。}
\studyroute{这里组织验证方案，略去具体实验不确定度、误差估计器和项目级数据管理；代数求解与计算架构见第 0 卷。}{verification and validation; grid convergence; temporal convergence; model error; reduction error; quantity of interest}{为一个波动或输运问题定义目标量，分别改变模型、网格、次数、步长与代数容差，建立能够区分误差来源的实验矩阵。}
\section{方法选择与组合策略}
几何固定且壁面精度要求高时，贴体网格便于直接解析边界层。

对于复杂运动几何，若主要目标位于壁面剪切应力之外，浸没边界法可简化网格组织。

若局部尺度变化极大，可使用四叉树或八叉树自适应网格。

若边界运动幅度较小，可使用动网格。

若多个刚体发生大幅相对运动，可使用重叠网格。

实际工业软件常组合使用：
\[
\text{贴体边界层}
+
\text{自适应远场},
\]
\[
\text{八叉树背景网格}
+
\text{浸没边界},
\]
\[
\text{重叠网格}
+
\text{局部动网格}.
\]

\begin{askai}
AI 展开问题：
\begin{quote}
给定一个带旋转叶片和远场声波传播的计算问题，请比较贴体网格、动网格、重叠网格和自适应远场网格各自解决什么子问题。请按各子区域的需求给出组合网格策略。
\end{quote}

进一步展开：
\begin{quote}
把“几何固定性、运动幅度、边界层精度、局部尺度变化、稳定步长”做成一个方法选择表。
\end{quote}
\end{askai}
\section{完整计算工作流}
一次可靠计算通常按以下链条进行：

\begin{enumerate}
  \item 用 \(\mathrm{Kn}\)、\(\mathrm{Re}\)、\(\mathrm{Ma}\) 等尺度比选择模型；
  \item 确定守恒变量、材料模型和边界条件；
  \item 选择与几何和局部尺度匹配的网格策略；
  \item 选择满足守恒、稳定性和波动精度要求的离散格式；
  \item 根据显式或隐式推进形成代数系统；
  \item 记录代数容差与目标量对残差的敏感性，相关求解算法见第 0 卷；
  \item 分别监控模型误差、网格误差、离散误差、线性残差和非线性残差。
\end{enumerate}

残差衡量代数方程的满足程度；模型有效性与离散精度还需通过基准解、网格及时间收敛实验和物理数据验证。

\section{五类误差与验证证据}\index{验证与确认}
\begin{center}\small
\begin{tabularx}{\textwidth}{@{}l>{\raggedright\arraybackslash}X@{}}
\toprule
误差层次 & 可执行检查\\
\midrule
物理模型 & 与实验、测量或适用范围内的高保真模型比较，报告参数和数据不确定度。\\
约化与闭合 & 改变小参数，比较完整与约化模型；检验未解矩和尺度项的量级。\\
空间离散 & 固定充分小的步长与容差，分别细化 $h$、增加 $p$，比较目标量和预期阶数。\\
时间离散 & 固定充分解析的空间，减半 $\Delta t$，检查时间阶数与相位累积。\\
代数求解 & 收紧残差容差，观察目标量变化是否已小于空间和时间误差。\\
\bottomrule
\end{tabularx}
\end{center}
几何近似、速度域截断、随机抽样和边界模型还应在相应层次单列。
代数残差需通过算子稳定性或目标量敏感性与解误差联系。

\section{守恒、稳定与收敛的实验}
在周期或封闭无源系统中，记录总质量、动量与能量；
有物理源和边界时，记录输入输出及耗散的平衡缺口。
能量估计或离散熵给出允许的演化范围，波动问题还需测量相位与群速度。
电磁和 MHD 问题同时记录离散散度。

光滑解析解可检验预期阶数，制造解可指定源项和边界来覆盖全部算子。
间断、壁面层和多尺度问题需采用适合其正则性和目标量的误差指标。
数值验证考察给定模型的求解误差；物理确认考察模型对目标现象的描述能力。

\conceptbridge{残差、守恒和网格收敛都通过后，模型还需什么证据？}
{这些检查针对特定离散和模型的数学表现。物理确认还需目标条件下的测量或适当基准，
并评估本构、边界和参数的不确定度。各层证据共同限定计算结论的可信范围。}
{数值验证 $\longrightarrow$ 模型与数据比较 $\longrightarrow$ 目标量的可信范围。}

\aiexercise{为周期 Alfv\'en 波或低 Mach 剪切波建立误差表，分别列出网格、次数、步长、容差、守恒缺口和目标量误差。设计能够识别时间误差掩盖空间收敛的对照组。}
{将目标换成壁面热流或湍流输运，列出新增的碰撞、本构、边界和统计假设，以及可以独立检验这些假设的数据。}

\section{全书的依赖链}
\[
\text{Liouville 与统计表示}
\longrightarrow\text{动理学与连续模型}
\longrightarrow\text{尺度约化与闭合}
\]
\[
\longrightarrow\text{函数空间与弱形式}
\longrightarrow\text{残差、通量与高阶离散}
\longrightarrow\text{网格与验证}.
\]
保留的状态决定所需闭合，连续算子与目标结构决定适合的离散约束，
目标量与验证证据再决定计算分辨率和方法组合。

\backmatter
\cleardoublepage
\phantomsection
\addcontentsline{toc}{chapter}{参考文献}
\begin{thebibliography}{99}
\raggedright
\bibitem{evans}
L. C. Evans, \emph{Partial Differential Equations}, 2nd ed., American Mathematical Society, 2010.
弱导数、Sobolev 空间、能量估计与椭圆问题的进一步阅读。
\bibitem{leveque}
R. J. LeVeque, \emph{Finite Volume Methods for Hyperbolic Problems}, Cambridge University Press, 2002.
守恒律、Riemann 问题与高分辨率通量。
\bibitem{hesthaven}
J. S. Hesthaven and T. Warburton, \emph{Nodal Discontinuous Galerkin Methods}, Springer, 2008.
局部基函数、节点表示、能量与谱分析。
\bibitem{jiang1996}
G.-S. Jiang and C.-W. Shu,
\href{https://doi.org/10.1006/jcph.1996.0130}{\emph{Efficient Implementation of Weighted ENO Schemes}},
Journal of Computational Physics, 126 (1996), 202--228.
\bibitem{arnold2002}
D. N. Arnold, F. Brezzi, B. Cockburn, and L. D. Marini,
\href{https://doi.org/10.1137/S0036142901384162}{\emph{Unified Analysis of Discontinuous Galerkin Methods for Elliptic Problems}},
SIAM Journal on Numerical Analysis, 39 (2002), 1749--1779.
\bibitem{huynh2007}
H. T. Huynh,
\href{https://doi.org/10.2514/6.2007-4079}{\emph{A Flux Reconstruction Approach to High-Order Schemes Including Discontinuous Galerkin Methods}},
AIAA Paper 2007-4079, 2007.
\bibitem{vincent2011}
P. E. Vincent, P. Castonguay, and A. Jameson,
\href{https://doi.org/10.1007/s10915-010-9420-z}{\emph{A New Class of High-Order Energy Stable Flux Reconstruction Schemes}},
Journal of Scientific Computing, 47 (2011), 50--72.
\bibitem{ranocha2016}
H. Ranocha, P. \"Offner, and T. Sonar,
\href{https://arxiv.org/abs/1511.02052}{\emph{Summation-by-parts operators for correction procedure via reconstruction}},
2016, arXiv:1511.02052.
\bibitem{kopriva1996}
D. A. Kopriva,
\href{https://doi.org/10.1006/jcph.1996.0225}{\emph{A Conservative Staggered-Grid Chebyshev Multidomain Method for Compressible Flows. II. A Semi-Structured Method}},
Journal of Computational Physics, 128 (1996), 475--488.
\bibitem{dumbser2014}
M. Dumbser,
\href{https://arxiv.org/abs/1402.6897}{\emph{Arbitrary-Lagrangian-Eulerian ADER-WENO Finite Volume Schemes with Time-Accurate Local Time Stepping for Hyperbolic Conservation Laws}},
2014, doi:10.1016/j.cma.2014.07.019.
\bibitem{owren1992}
B. Owren and M. Zennaro,
\href{https://doi.org/10.1137/0913084}{\emph{Derivation of Efficient, Continuous, Explicit Runge--Kutta Methods}}, 1992.
\bibitem{altmann2010}
C. Altmann, G. Gassner, and C.-D. Munz,
\href{https://congress2.cimne.com/eccomas/proceedings/cfd2010/papers/01008.pdf}{\emph{An Explicit Space-Time Adaptive Discontinuous Galerkin Scheme}},
ECCOMAS CFD, 2010.
\bibitem{ascher1997}
U. M. Ascher, S. J. Ruuth, and R. J. Spiteri,
\href{https://doi.org/10.1016/S0168-9274(97)00056-1}{\emph{Implicit-explicit Runge--Kutta methods for time-dependent partial differential equations}},
Applied Numerical Mathematics, 25 (1997), 151--167.
\bibitem{qian1992}
Y. H. Qian, D. d'Humi\`eres, and P. Lallemand,
\href{https://doi.org/10.1209/0295-5075/17/6/001}{\emph{Lattice BGK Models for Navier-Stokes Equation}},
Europhysics Letters, 17 (1992), 479--484.
\bibitem{dhumieres2002}
D. d'Humi\`eres, I. Ginzburg, M. Krafczyk, P. Lallemand, and L.-S. Luo,
\href{https://doi.org/10.1098/rsta.2001.0955}{\emph{Multiple-relaxation-time lattice Boltzmann models in three dimensions}},
Philosophical Transactions of the Royal Society A, 360 (2002), 437--451.
\bibitem{geier2015}
M. Geier, M. Sch\"onherr, A. Pasquali, and M. Krafczyk,
\href{https://doi.org/10.1016/j.camwa.2015.05.001}{\emph{The cumulant lattice Boltzmann equation in three dimensions: Theory and validation}},
Computers \& Mathematics with Applications, 70 (2015), 507--547.
\bibitem{karlin1999}
I. V. Karlin, A. Ferrante, and H. C. \"Ottinger,
\href{https://doi.org/10.1209/epl/i1999-00370-1}{\emph{Perfect entropy functions of the Lattice Boltzmann method}},
Europhysics Letters, 47 (1999), 182--188.
\bibitem{grad1949}
H. Grad,
\href{https://doi.org/10.1002/cpa.3160020403}{\emph{On the kinetic theory of rarefied gases}},
Communications on Pure and Applied Mathematics, 2 (1949), 331--407.
\bibitem{holway1966}
L. H. Holway,
\href{https://doi.org/10.1063/1.1761920}{\emph{New Statistical Models for Kinetic Theory: Methods of Construction}},
Physics of Fluids, 9 (1966), 1658--1673.
\bibitem{bird1994}
G. A. Bird,
\href{https://doi.org/10.1093/oso/9780198561958.001.0001}{\emph{Molecular Gas Dynamics and the Direct Simulation of Gas Flows}},
Clarendon Press, 1994.
\bibitem{cgl1956}
G. F. Chew, M. L. Goldberger, and F. E. Low,
\href{https://static.ias.edu/pitp/2016/sites/pitp/files/chew_goldberger_low_1956-1.pdf}{\emph{The Boltzmann equation and the one-fluid hydromagnetic equations in the absence of particle collisions}},
Proceedings of the Royal Society A, 236 (1956), 112--118.
\bibitem{brizard2007}
A. J. Brizard and T. S. Hahm,
\href{https://doi.org/10.1103/RevModPhys.79.421}{\emph{Foundations of nonlinear gyrokinetic theory}},
Reviews of Modern Physics, 79 (2007), 421--468.
\href{https://bp-pub.pppl.gov/pub_report/2006/PPPL-4153.pdf}{PPPL 预印本}。
\bibitem{abel2013}
I. G. Abel et al.,
\href{https://arxiv.org/abs/1209.4782}{\emph{Multiscale Gyrokinetics for Rotating Tokamak Plasmas: Fluctuations, Transport and Energy Flows}},
Reports on Progress in Physics, 76 (2013), 116201.
\end{thebibliography}

\clearpage
\phantomsection
\addcontentsline{toc}{chapter}{术语索引}
\printindex
\end{document}
